\documentclass[11pt,letterpaper]{article}
\usepackage[T1]{fontenc}
\usepackage[utf8]{inputenc}
\usepackage{lmodern}
\usepackage[margin=1in,headheight=15pt]{geometry}
\usepackage{amsmath,amssymb,amsthm,mathtools,mathrsfs}
\usepackage{microtype,booktabs,longtable,tabularx,array,enumitem}
\usepackage[dvipsnames]{xcolor}
\usepackage{fancyhdr,aliascnt}
\usepackage[colorlinks=true,linkcolor=MidnightBlue,citecolor=MidnightBlue,urlcolor=MidnightBlue]{hyperref}
\usepackage[nameinlink,noabbrev]{cleveref}
\hypersetup{pdftitle={Surcomplex Polynomial Algebra: Factorization, Root Geometry, Multiscale Stability, and Residue Duality},pdfauthor={Research exposition prepared with ChatGPT},pdfsubject={Finite polynomial algebra over surreal and Hahn fields}}
\pagestyle{fancy}\fancyhf{}
\fancyhead[L]{\small Surcomplex polynomial algebra}
\fancyhead[R]{\small Roots, scales, and duality}
\fancyfoot[C]{\thepage}
\setlength{\parskip}{0.3em}
\setlength{\parindent}{1.25em}
\setlist{leftmargin=2em,itemsep=0.15em,topsep=0.3em}
\setcounter{tocdepth}{2}
\numberwithin{equation}{section}
\newtheorem{theorem}{Theorem}[section]
\newaliascnt{proposition}{theorem}
\newtheorem{proposition}[proposition]{Proposition}
\aliascntresetthe{proposition}
\newaliascnt{lemma}{theorem}
\newtheorem{lemma}[lemma]{Lemma}
\aliascntresetthe{lemma}
\newaliascnt{corollary}{theorem}
\newtheorem{corollary}[corollary]{Corollary}
\aliascntresetthe{corollary}
\theoremstyle{definition}
\newaliascnt{definition}{theorem}
\newtheorem{definition}[definition]{Definition}
\aliascntresetthe{definition}
\newaliascnt{example}{theorem}
\newtheorem{example}[example]{Example}
\aliascntresetthe{example}
\theoremstyle{remark}
\newaliascnt{remark}{theorem}
\newtheorem{remark}[remark]{Remark}
\aliascntresetthe{remark}
\crefname{theorem}{theorem}{theorems}
\Crefname{theorem}{Theorem}{Theorems}
\crefname{proposition}{proposition}{propositions}
\Crefname{proposition}{Proposition}{Propositions}
\crefname{lemma}{lemma}{lemmas}
\Crefname{lemma}{Lemma}{Lemmas}
\crefname{corollary}{corollary}{corollaries}
\Crefname{corollary}{Corollary}{Corollaries}
\crefname{example}{example}{examples}
\Crefname{example}{Example}{Examples}
\crefname{remark}{remark}{remarks}
\Crefname{remark}{Remark}{Remarks}
\crefname{definition}{definition}{definitions}
\Crefname{definition}{Definition}{Definitions}
\newcommand{\No}{\mathbf{No}}
\newcommand{\SC}{\mathbf{SC}}
\newcommand{\C}{\mathbb C}
\newcommand{\R}{\mathbb R}
\newcommand{\Q}{\mathbb Q}
\newcommand{\N}{\mathbb N}
\newcommand{\Z}{\mathbb Z}
\newcommand{\OO}{\mathcal O}
\newcommand{\HH}{\mathcal H_\Gamma}
\newcommand{\Ov}{\mathcal O_v}
\newcommand{\mv}{\mathfrak m_v}
\newcommand{\sh}{^{\sharp}}
\newcommand{\supp}{\operatorname{supp}}
\newcommand{\st}{\operatorname{st}}
\newcommand{\red}{\operatorname{red}}
\newcommand{\lc}{\operatorname{lc}}
\newcommand{\ord}{\operatorname{ord}}
\newcommand{\Res}{\operatorname{Res}}
\newcommand{\Disc}{\operatorname{Disc}}
\newcommand{\Tr}{\operatorname{Tr}}
\newcommand{\Nm}{\operatorname{N}}
\newcommand{\remp}{\operatorname{rem}}
\newcommand{\conv}{\operatorname{conv}}
\newcommand{\Spec}{\operatorname{Spec}}
\newcommand{\id}{\operatorname{id}}
\newcommand{\abs}[1]{\left|#1\right|}
\newcommand{\Bge}[2]{B_{\ge}(#1,#2)}
\newcommand{\Bgt}[2]{B_{>}(#1,#2)}
\newcolumntype{L}{>{\raggedright\arraybackslash}X}
\newcolumntype{P}[1]{>{\raggedright\arraybackslash}p{#1}}
\begin{document}
\hypersetup{pageanchor=false}
\begin{titlepage}
\centering
\vspace*{0.7cm}
{\large\scshape A finite-algebra continuation of the supplied manuscripts\par}
\vspace{0.9cm}
{\Huge\bfseries Surcomplex\par Polynomial Algebra\par}
\vspace{0.6cm}
{\LARGE Factorization, Root Geometry,\par Multiscale Stability,\par and Residue Duality\par}
\vspace{0.85cm}
{\large September 21, 2026\par}
\vspace{0.7cm}
\begin{minipage}{0.94\textwidth}
\begin{abstract}
We develop finite-degree polynomial algebra over the surcomplex class field
$\SC=\No[i]$, using set-sized algebraically closed Hahn fields for every
individual construction. Two complementary geometries are essential:
order-valued modulus geometry and natural-valuation geometry. In the first,
we prove Gauss--Lucas directly over the surreal real field, obtain root
bounds, and justify a polynomial disk version of Rouch\'e by real-closed-field
transfer. In the second, a single normalized reduction polynomial counts
roots in closed balls, open balls, and individual residue classes at any
surreal center and scale. It yields Newton polygons, a valuation form of
Rouch\'e, and an exact critical-point count for every ball containing roots.

We derive an explicit perturbation threshold preserving the complete
pairwise valuation geometry of a squarefree root configuration. For monic
polynomials with finite roots, coefficient precision strictly beyond the
valuation of the discriminant is sufficient; a quadratic example shows why
strictness cannot be dropped. We also prove exact coefficient-uncertainty
root loci, support-controlled coprime factor lifting with holomorphic
parameters, and a finite cluster formula for discriminant valuation.

The algebraic part includes Hermite interpolation, resultants, Newton sums,
a universal residue pairing with an explicit dual basis, collision-stable
trace formulas, finite polynomial maps and ramification, and a class-safe
Nullstellensatz for finite polynomial systems. All infinite constructions
are justified by Hahn support, not by convergence of ordinary partial sums
in the fine surreal topology.
\end{abstract}
\end{minipage}
\vfill
\begin{minipage}{0.94\textwidth}\small
\textbf{Status and attribution.}
This is a theorem-driven research exposition, not a claim to resolve a named
open conjecture. Classical algebraic, real-closed, and valued-field
principles are identified as such. The supplied manuscripts are credited
for their Hahn-coherent framework and analytic antecedents. The particular
organization and derived estimates are not a certified priority claim.
The proofs have not been independently refereed or machine checked.
\end{minipage}
\end{titlepage}
\hypersetup{pageanchor=true}
\pagenumbering{roman}
\tableofcontents
\clearpage
\pagenumbering{arabic}

\section{The subject, the sources, and the main results}\label{sec:introduction}

\subsection{Why polynomial algebra deserves a separate treatment}
The definition of a surcomplex number is simple: $x+iy$ with $x,y\in\No$
and $i^2=-1$. The algebraic closure of this class field is established
surreal theory, not a new fundamental theorem of algebra
\cite{Gonshor,Poonen}. What remains worth developing is a usable polynomial
calculus that distinguishes finite algebra, ordered geometry, valuation
geometry, and support-controlled deformation. These distinctions become
visible when roots are infinite, are infinitesimally close, or collide.

Throughout, a polynomial has \emph{ordinary finite degree}. Its coefficients
may have arbitrary set-sized Hahn supports. We do not use ordinal degrees,
proper-class sums, or a polynomial-like series of unbounded degree under
the name ``polynomial.'' All theorems about an individual finite root
configuration are proved in a set-sized field containing its data and then
interpreted in $\SC$.

Three supplied manuscripts form the starting point. The first,
\emph{Surcomplex Analysis} \cite{SourceAnalysis}, separates fine-local
analytic germs, Hahn-coherent functions on thickened ordinary domains,
and coherence at every surreal scale. The second, \emph{Hartogs Coherence,
Finite Maps, and Residue Duality} \cite{SourceResearch}, develops
fixed-domain parameter division, finite intersection algebras, and
collision-stable residue functionals. The third, \emph{Global Divisors and
Cohomological Obstructions} \cite{SourceGlobal}, identifies support
obstructions for infinitely many prescribed clusters and interpolation
conditions. We retain their terminology and summation convention.

The emphasis here is different. We work with actual finite polynomials at
\emph{all} centers and scales, prove geometric statements requiring only
finite algebra, and make precise how coefficient precision controls a
complete root configuration. No global analytic assertion from the sources
is silently enlarged beyond its stated category.

\subsection{A map of the results}
The central statements can be read in three groups.

\paragraph{Root geometry.}
The Gauss--Lucas theorem holds with convex combinations in the real surreal
field (\cref{thm:gausslucas}). Polynomial Rouch\'e holds on genuine
order-modulus circles, including infinite and infinitesimal circles
(\cref{thm:orderrouche}). At any center $a$ and value-group scale $\delta$,
normalize $P(a+t^\delta U)$ to have least coefficient valuation zero and
reduce its coefficients to $\C$. The resulting nonzero polynomial $R(U)$
has degree equal to the number of roots in $\Bge{a}{\delta}$ and order at
zero equal to the number in $\Bgt{a}{\delta}$
(\cref{thm:reductioncount}). Its other root multiplicities identify the
finer residue clusters.

\paragraph{Critical points and stability.}
A valuation ball containing $m\ge1$ roots of $P$, counted with multiplicity,
contains exactly $m-1$ roots of $P'$ (\cref{thm:criticalballs}). For a
monic polynomial with distinct finite roots $\alpha_1,\ldots,\alpha_n$, put
\begin{equation}\label{eq:introthreshold}
 \sigma_i=\sum_{j\ne i}v(\alpha_i-\alpha_j),\qquad
 \delta_i=\max_{j\ne i}v(\alpha_i-\alpha_j),\qquad
 T(P)=\max_i(\sigma_i+\delta_i).
\end{equation}
If another monic degree-$n$ polynomial has coefficient errors of valuation
strictly greater than $T(P)$, its roots have a unique matching with the
$\alpha_i$ preserving every pairwise distance valuation. The more readily
computed condition
\[
 \min_jv(q_j-p_j)>v(\Disc P)
\]
is sufficient (\cref{thm:stability,cor:discprecision}). This conclusion is
stronger than merely preserving the number of roots.

\paragraph{Finite algebra and support.}
We prove finite Hermite interpolation and Chinese remainders, exact root
loci under bounded coefficient errors, and coprime factor lifting whose
support lies in the positive monoid generated by the perturbation
(\cref{thm:CRT,thm:pseudozeros,thm:hensel}). The residue functional
\[
 \lambda_P(h)=[X^{n-1}]\remp_P h
\]
has a perfect pairing for every monic $P$, even over a non-Noetherian
coefficient ring and through repeated-root collisions. An explicit dual
basis gives $\Tr(m_h)=\lambda_P(P'h)$ without an inverse discriminant
(\cref{thm:residuepairing,thm:trace}). Finite polynomial maps, ramification,
and a finite-system Nullstellensatz complete the algebraic picture.

\subsection{What is classical, and what is being added}
Euclidean division, finite Chinese remainders, resultants, the
Nullstellensatz, and residue duality are classical algebra. Gauss--Lucas is
classical complex polynomial geometry, and its proof below works over any
real closed field. Real-closed-field transfer and non-Archimedean root
geometry also have extensive precedents \cite{CohenMahboubi,Eisermann,Faber}.
It would be misleading to label their surcomplex instances wholly new.

The contribution of this article is an explicit, mutually compatible
package: arbitrary-rank Hahn normalization, exact ball and residue counts,
critical-point allocation, support bounds, and a quantitative
cluster-preservation theorem, together with proofs and sharp examples.
The particular formulations are presented for scrutiny, not as certified
first occurrences in the literature. \Cref{sec:scope} and the supplementary
literature audit distinguish inherited facts from the derivations here.

\section{Foundations: finite degree, Hahn fields, and two geometries}\label{sec:foundations}

\subsection{Set-sized fields inside the surcomplex class}
We use class set theory, for example NBG with global choice, in the same
way as the supplied manuscripts. Each individual support, recursion, list
of coefficients, and algebraic certificate is a set. Conway normal forms
give, for a set-sized ordered abelian subgroup $\Gamma\subseteq\No$,
\begin{equation}\label{eq:workspace}
 K_\Gamma=\C((t^\Gamma))\subseteq\SC,\qquad
 t^\gamma=\omega^{-\gamma}.
\end{equation}
The notation $t^\gamma$ means a Conway monomial. An element is a sum
$\sum_{\gamma\in S}c_\gamma t^\gamma$ with $S$ well ordered in the increasing
order and $c_\gamma\in\C$. For a nonzero element define
\[
 v(x)=\min\supp(x),\qquad \lc(x)=c_{v(x)},\qquad v(0)=+\infty.
\]
Then $v(xy)=v(x)+v(y)$ and $v(x+y)\ge\min(v(x),v(y))$.

\begin{proposition}[A field for every set of data]\label{prop:workspace}
Every set of surcomplex constants is contained in a field $K_\Gamma$ with
$\Gamma$ divisible and set-sized. Such a field is algebraically closed.
Its fixed field under complex conjugation is the real closed ordered
field $R_\Gamma=\R((t^\Gamma))$, and $K_\Gamma=R_\Gamma[i]$.
\end{proposition}
\begin{proof}
Take the union of the normal-form supports of the given constants,
generate its additive subgroup in $\No$, and take the divisible hull.
These are set operations. The algebraic closure of the resulting Hahn
field is the classical theorem for divisible value group and algebraically
closed coefficient field; see Poonen \cite[Corollary~4]{Poonen}.

Real and imaginary parts of a Hahn series lie in $R_\Gamma$, so
$K_\Gamma=R_\Gamma[i]$. The leading-coefficient order makes $R_\Gamma$ an
ordered field. A positive element has a square root in $K_\Gamma$; writing
that root as $x+iy$ shows $xy=0$, and positivity excludes a purely
imaginary square root. Every odd-degree polynomial over $R_\Gamma$ has a
real root, because its nonreal roots over $K_\Gamma$ occur in conjugate
pairs. The standard characterization of real closed ordered fields now
applies.
\end{proof}

A nonzero polynomial over $K_\Gamma$ already splits there. Enlarging the
workspace does not add roots: a product of its linear factors cannot
vanish at any additional element of a field extension. The same observation
preserves root multiplicities. We shall enlarge workspaces to include
centers, scales, and finitely many parameters without further comment.

\begin{remark}[What is not a set ring]
The notation $\SC[X]$ is convenient for the class of finite polynomials,
but it does not authorize an application of Zorn's lemma to a proper-class
family of ideals. Euclidean calculations concern finitely many
polynomials. Maximal ideals, finite-dimensional algebras, and the
Nullstellensatz below are first handled in the set ring $K_\Gamma[X]$ or
$K_\Gamma[X_1,\ldots,X_d]$.
\end{remark}

\subsection{Modulus, valuation, and reduction}
Conjugation and modulus are
\[
 \overline{x+iy}=x-iy,\qquad |x+iy|=\sqrt{x^2+y^2}.
\]
The positive square root exists in the real closed field. The usual
algebraic proofs give $|zw|=|z||w|$ and $|z+w|\le|z|+|w|$, and normal forms
give $v(|z|)=v(z)$. Let
\[
 \Ov=\{z:v(z)\ge0\},\qquad \mv=\{z:v(z)>0\}.
\]
These are the finite elements and infinitesimals, respectively. Reduction
$\red:\Ov\to\C$ extracts the exponent-zero coefficient; it is the
standard-part map. We reserve an overline for complex conjugation, not for
this reduction. Every finite $z$ is uniquely $\red(z)+\varepsilon$ with
$\varepsilon\in\mv$.

For $a\in K_\Gamma$ and $\delta\in\Gamma$ define
\begin{equation}\label{eq:valuationballs}
 \Bge{a}{\delta}=\{z:v(z-a)\ge\delta\}=a+t^\delta\Ov,
 \qquad
 \Bgt{a}{\delta}=\{z:v(z-a)>\delta\}=a+t^\delta\mv.
\end{equation}
By contrast, an order-modulus disk is $\{z:|z-a|<r\}$ for a positive surreal
$r$. These are not interchangeable definitions. For instance,
$\Bge{0}{0}$ contains every finite complex number, not just the modulus unit
disk. The sets in \eqref{eq:valuationballs} are called closed and open
\emph{valuation} balls; this terminology does not assert a special role for
them in the full fine class topology.

If $v(x)>v(y)$ with $y\ne0$, then $x/y$ is infinitesimal, so $|x|<|y|/n$
for every positive integer $n$. Valuation geometry forgets ordinary
nonzero complex factors while recording every surreal scale. This is why
its perturbation criteria differ from modulus criteria.

\subsection{The support lemma and the absence of a convergence shortcut}
A family of Hahn series is \emph{strongly summable} if the union of its
supports is well ordered and only finitely many members contribute at each
exponent. For a positive well-ordered set $E\subseteq\Gamma_{>0}$ write
\[
 E^*=\{e_1+\cdots+e_k:k\ge0,\ e_j\in E\},
\]
including the empty sum $0$.

\begin{lemma}[Hahn--Neumann support calculus]\label{lem:support}
If $S,T$ are well ordered in an ordered abelian group, then $S+T$ is well
ordered and each element has finitely many representations $s+t$.
If $E$ is well ordered and positive, $E^*$ is well ordered and each exponent
has finitely many representations by nondecreasing finite words in $E$.
Consequently an expression indexed by one exponent of $S$ and a finite
word in $E$ is strongly summable when each such word contributes only
finitely many terms.
\end{lemma}
\begin{proof}
The first two assertions are the classical Hahn--Neumann lemma
\cite{Neumann}; the same formulation is used in all three supplied
manuscripts. For the last assertion, first decompose an output exponent as
$s+\eta$ with $s\in S$, $\eta\in E^*$. There are finitely many such pairs.
For each $\eta$ there are finitely many nondecreasing words, and each finite
word has finitely many permutations. These are exactly the two
requirements for strong summability.
\end{proof}

The lemma justifies $(1+h)^{-1}=\sum_{k\ge0}(-h)^k$ for positive-support
$h$. It does \emph{not} say that these partial sums converge in the fine
surreal topology. Nor does it say that $k\delta$ eventually exceeds every
positive value-group element. The latter is false in higher rank. All
recursive constructions below are verified coefficient by coefficient
using finite support words.

\subsection{Finite polynomial data inside the coherent framework}
For an ordinary connected domain $U\subseteq\C^d$ the supplied category is
\[
 \HH(U)=\OO(U)((t^\Gamma)).
\]
A datum $F=\sum f_\gamma t^\gamma$ is evaluated at a finite displaced point
$c+\varepsilon$ by Taylor expanding each ordinary holomorphic coefficient:
\begin{equation}\label{eq:coherentevaluation}
 F\sh(c+\varepsilon)=
 \sum_{\gamma}\sum_{\alpha\in\N^d}
 \frac{\partial^\alpha f_\gamma(c)}{\alpha!}
 t^\gamma\varepsilon^\alpha.
\end{equation}
The support lies in $\supp(F)+(\bigcup_j\supp\varepsilon_j)^*$, and
\cref{lem:support} verifies finite multiplicity. This is inherited from
\cite{SourceAnalysis,SourceResearch}.

In this paper $\mathcal H_{\Gamma,\ge0}$ and
$\mathcal H_{\Gamma,>0}$ always mean nonnegative and strictly positive
support. This avoids a notational mismatch between the supplied papers,
where a superscript or subscript $+$ is used for different conventions.

A polynomial with coefficients in $\HH(U)$ has a uniform finite degree
bound in $X$. Conversely, if every coefficient polynomial in a Hahn datum
has degree at most the same finite $n$, collecting the $n+1$ powers of $X$
gives such a polynomial. Without that bound this fails. For example
\begin{equation}\label{eq:geometricwarning}
 \sum_{k\ge0}t^kX^k
\end{equation}
is a legitimate coherent datum on the thickening of every ordinary domain,
and equals $(1-tX)^{-1}$ there. It is not a finite polynomial and is not
a global pole-free function of all surreal scales: the rational function
has a pole at $X=t^{-1}$. This elementary example separates the subject of
this article from unrestricted Hahn-coherent entire data.

\section{Factorization, finite interpolation, and symmetric identities}\label{sec:finitealgebra}
Throughout this section $K$ is an algebraically closed workspace as in
\cref{prop:workspace}. All statements therefore apply to surcomplex
polynomials.

\subsection{Euclidean algebra and multiplicity}
\begin{theorem}[Finite polynomial algebra]\label{thm:euclidean}
For $A,B\in K[X]$, $B\ne0$, there are unique $Q,R\in K[X]$ with
$A=BQ+R$ and $\deg R<\deg B$. Every pair has a monic greatest common
divisor $D$, and $D=UA+VB$ for some $U,V\in K[X]$ when the pair is not
identically zero. Every nonconstant polynomial has a factorization
\begin{equation}\label{eq:factorization}
 P(X)=a_n\prod_{i=1}^s(X-\alpha_i)^{m_i},
 \qquad \alpha_i\ne\alpha_j\ (i\ne j),\qquad \sum_i m_i=n.
\end{equation}
The factors and multiplicities are unique up to their order. The
multiplicity at $\alpha$ is the least $j$ for which $P^{(j)}(\alpha)\ne0$.
\end{theorem}
\begin{proof}
Successively cancel the highest term of $A$ by a multiple of $B$. At each
step the degree strictly decreases, so the procedure terminates after
finitely many steps. Subtracting two putative divisions proves uniqueness
by degree. Applying the same divisions to successive remainders gives the
Euclidean algorithm; reversing its finite identities gives B\'ezout
coefficients. Algebraic closure supplies a root of every nonconstant
factor, and division by $X-\alpha$ reduces degree, proving
\eqref{eq:factorization}. The factor theorem and cancellation in a domain
give uniqueness. Finally the finite Taylor identity
\[
 P(\alpha+U)=\sum_{j=0}^n\frac{P^{(j)}(\alpha)}{j!}U^j
\]
identifies the multiplicity, since the characteristic is zero.
\end{proof}

In particular $P$ is squarefree exactly when $\gcd(P,P')=1$. A repeated
root of multiplicity $m$ has multiplicity exactly $m-1$ in $P'$. The field
has characteristic zero, so no inseparable exception is hidden in these
statements.

\subsection{Chinese remainders and Hermite data}
Assume $P$ is monic and use \eqref{eq:factorization}. Let
$A_i=(X-\alpha_i)^{m_i}$ and $P_i=P/A_i$. The inverse of $P_i$ modulo $A_i$
is the finite Taylor polynomial at $\alpha_i$ of the rational function
$1/P_i$:
\begin{equation}\label{eq:inversejet}
 C_i(X)=\sum_{j=0}^{m_i-1}\frac1{j!}
 \left(\frac{d^j}{dX^j}\frac1{P_i(X)}\right)_{X=\alpha_i}
 (X-\alpha_i)^j.
\end{equation}
Its denominators are nonzero because the distinct $\alpha_i$ are separated
as field elements, even when their differences are infinitesimal.

\begin{theorem}[Finite Hermite interpolation and CRT]\label{thm:CRT}
The map
\begin{equation}\label{eq:finiteCRT}
 K[X]/(P)\longrightarrow\prod_{i=1}^s K[U]/(U^{m_i}),
 \qquad [H]\longmapsto[H(\alpha_i+U)]_i
\end{equation}
is an isomorphism. Arbitrary jets of orders less than $m_i$ at the
$\alpha_i$ have a unique polynomial interpolant of degree less than $n$.
The residue classes
\[
 E_i=[P_iC_i]\in K[X]/(P)
\]
are orthogonal idempotents with sum $1$ and give the identities of the
factors in \eqref{eq:finiteCRT}.
\end{theorem}
\begin{proof}
The finite Taylor identity gives $P_iC_i\equiv1\pmod{A_i}$. For $j\ne i$,
$P_iC_i\equiv0\pmod{A_j}$. Given target Taylor polynomials
$J_i\in K[U]_{<m_i}$, the remainder modulo $P$ of
\[
 \sum_i P_i(X)C_i(X)J_i(X-\alpha_i)
\]
has all the required jets. A polynomial in the kernel is divisible by
each $A_i$. These polynomials are pairwise coprime, so repeated use of
B\'ezout gives divisibility by their product $P$. This proves injectivity
and the degree-$<n$ uniqueness. The idempotent identities hold in every
factor and therefore in the quotient.
\end{proof}

When all roots are simple this becomes the Lagrange formula
\begin{equation}\label{eq:lagrange}
 H(X)=\sum_{i=1}^n H(\alpha_i)
 \frac{P(X)}{(X-\alpha_i)P'(\alpha_i)},\qquad \deg H<n.
\end{equation}
The denominators may be infinitely large after inversion, but the formula
is a finite algebraic identity. No separation lower bound in ordinary
real units is required.

\subsection{Vieta and Newton sums at arbitrary magnitude}
Write
\[
 P(X)=X^n+c_1X^{n-1}+\cdots+c_n,
 \qquad s_k=\sum_{i=1}^n\alpha_i^k,
\]
where roots are now listed with repetitions. Expanding the product gives
$c_j=(-1)^j e_j(\alpha_1,\ldots,\alpha_n)$. The power sums satisfy
\begin{align}
 s_k+c_1s_{k-1}+\cdots+c_{k-1}s_1+kc_k&=0 &&(1\le k\le n),\label{eq:newton1}\\
 s_k+c_1s_{k-1}+\cdots+c_ns_{k-n}&=0 &&(k>n).\label{eq:newton2}
\end{align}
Here a displayed sum is finite, even if its summands are infinite
surcomplex numbers that cancel.

For a proof, use the formal identity in the Laurent series field at
infinity,
\[
 \frac{P'(X)}{P(X)}=\sum_{i=1}^n\frac1{X-\alpha_i}
 =\sum_{k\ge0}s_kX^{-k-1},\qquad s_0=n.
\]
This is a formal expansion in the independent variable $X^{-1}$, not an
evaluation at a particular surreal $X$. Multiply by $P$ and compare its
finitely determined coefficient at $X^{n-k-1}$. This proves
\eqref{eq:newton1}--\eqref{eq:newton2}. The same identities will reappear as
traces of multiplication in the finite quotient algebra.

\section{Order-valued polynomial geometry}\label{sec:ordergeometry}

\subsection{Root bounds without Archimedean arguments}
\begin{proposition}[Two root bounds]\label{prop:rootbounds}
Let $P(X)=a_nX^n+\cdots+a_0$ with $a_n\ne0$ and $n\ge1$. Every root
$\alpha$ satisfies
\begin{equation}\label{eq:cauchybound}
 |\alpha|\le1+\max_{j<n}|a_j/a_n|.
\end{equation}
If at least one lower coefficient is nonzero, also
\begin{equation}\label{eq:scaledbound}
 |\alpha|<2\max_{j<n}|a_j/a_n|^{1/(n-j)}.
\end{equation}
All powers and inequalities are in the real closed surreal field.
\end{proposition}
\begin{proof}
Divide by $a_n$ and let $M=\max_{j<n}|a_j/a_n|$. If
$r=|\alpha|>1+M$, then the finite geometric identity gives
\[
 \sum_{j<n}|a_j/a_n|r^j
 \le M\frac{r^n-1}{r-1}<r^n,
\]
contradicting $P(\alpha)=0$. For the second bound let the maximum on the
right without the factor $2$ be $B>0$. If $r\ge2B$, then
\[
 \sum_{j<n}|a_j/a_n|r^{j-n}
 \le\sum_{k=1}^n(B/r)^k
 \le\sum_{k=1}^n2^{-k}<1,
\]
again contradicting the root equation. These are finite inequalities;
no real sequence or limit is involved.
\end{proof}

Applying \eqref{eq:cauchybound} to the reversed polynomial gives a lower
bound on nonzero roots when $a_0\ne0$. More importantly, translating and
scaling by a sufficiently large positive surreal turns every finite
polynomial problem into one whose roots are all finite. The precision
change under this operation is recorded in \cref{subsec:rescaledprecision}.

\subsection{Gauss--Lucas with surreal convex weights}
For a finite subset of $K=R_\Gamma[i]$, its convex hull means all finite
linear combinations with nonnegative coefficients in $R_\Gamma$ summing
to $1$. In the full class interpretation the coefficients may be real
surreals. Restricting them to ordinary real numbers would be a different,
and generally too small, hull.

\begin{theorem}[Surcomplex Gauss--Lucas]\label{thm:gausslucas}
Every zero of $P'$ lies in the convex hull of the zeros of a nonconstant
surcomplex polynomial $P$. More explicitly, if $P'(z)=0$ and $P(z)\ne0$,
then, with the distinct roots and multiplicities of \eqref{eq:factorization},
\begin{equation}\label{eq:barycentric}
 z=\frac{\displaystyle\sum_i
        \frac{m_i\alpha_i}{|z-\alpha_i|^2}}
       {\displaystyle\sum_i\frac{m_i}{|z-\alpha_i|^2}}.
\end{equation}
All weights in this formula are strictly positive real surreals.
\end{theorem}
\begin{proof}
At a point that is not a root, logarithmic differentiation gives
\[
 0=\frac{P'(z)}{P(z)}=\sum_i\frac{m_i}{z-\alpha_i}.
\]
Conjugate this equality and use
$1/\overline{z-\alpha_i}=(z-\alpha_i)/|z-\alpha_i|^2$.
The result is
$z\sum_i m_i/|z-\alpha_i|^2=\sum_i m_i\alpha_i/|z-\alpha_i|^2$.
The denominator is a finite sum of positive elements and hence is nonzero.
This proves \eqref{eq:barycentric}. A critical point that is already a
root is in the hull without this calculation.
\end{proof}

This is a direct algebraic proof of the classical theorem over a real
closed field, not a use of compactness or separating hyperplanes on the
surreal class. Ordinary and formally verified complex versions are part
of the prior literature \cite{Eberl}.

Iterating the theorem places every zero of $P^{(k)}$, for $1\le k<n$, in
the same hull. In particular, every convex order-modulus disk containing
all roots contains all critical points. The average of the $n-1$
critical points, listed with multiplicity, also equals the average of the
$n$ roots: comparison of the two highest coefficients gives
\[
 \frac1{n-1}\sum_{P'(c)=0}c=\frac1n\sum_{P(\alpha)=0}\alpha
 \qquad(n\ge2).
\]
Neither conclusion describes how critical points are distributed among
infinitesimal clusters. That more precise question is answered in
\cref{sec:critical}.

\subsection{Polynomial Rouch\'e by real-closed-field transfer}
We use one standard logical input. The theory of real closed fields has
quantifier elimination and is complete. Thus a first-order sentence in
the ordered-field language true in $\R$ is true in every real closed
field; see \cite{Tarski,CohenMahboubi}. For a fixed degree, a complex polynomial is
encoded by pairs of real coefficient variables. Equations and inequalities
in squared moduli are polynomial formulas in those variables.

\begin{theorem}[Rouch\'e on an order-modulus disk]\label{thm:orderrouche}
Let $P,Q\in\SC[X]$ be nonzero, $a\in\SC$, and $r\in\No_{>0}$. Suppose
\begin{equation}\label{eq:orderrouchehyp}
 |Q(z)-P(z)|<|P(z)|\qquad\text{for every }z\in\SC
 \text{ with }|z-a|=r.
\end{equation}
Then $P$ and $Q$ have the same number of roots in $|z-a|<r$, counted with
multiplicity, and neither has a root on the boundary.
\end{theorem}
\begin{proof}
Fix the two finite degrees. Boundary condition \eqref{eq:orderrouchehyp}
is a first-order formula: quantify the two real coordinates of $z$, write
the circle as $(\Re z-\Re a)^2+(\Im z-\Im a)^2=r^2$, and square both
nonnegative sides of the inequality.

The conclusion is also first order at these fixed degrees. Introduce a
list of roots for each polynomial, with repetition, require that its
linear-factor product equals the polynomial coefficient by coefficient,
and count inside roots by a finite disjunction over subsets of the list.
This describes multiplicities even when several list entries coincide.
There is no quantification over an unbounded natural degree inside the
ordered-field language.

The resulting universally quantified implication is true over $\R$ by
ordinary polynomial Rouch\'e. It therefore holds over every real closed
field. Choose $R_\Gamma$ containing the real and imaginary parts of the
coefficients, $a$, and $r$. The hypothesis holds in that field because it
holds in $\SC$. All roots already lie in $R_\Gamma[i]$, so its conclusion
is precisely the desired count in $\SC$. The strict inequality excludes
zeros of $P$ on the circle; if $Q$ vanished there it would give equality of
the two moduli, so $Q$ has no boundary root either.
\end{proof}

This is a proof by transfer from a classical theorem, and is identified as
such. It does not define an integral along a fine-continuous real-parameter
surreal circle. Constructive real-algebraic winding-number methods offer
another route to polynomial root counting over real closed fields
\cite{Eisermann}; we do not need that machinery for the stated implication.

\begin{corollary}[A coefficient-dominance root count]\label{cor:pellet}
Write $P(a+Z)=\sum_{j=0}^n p_jZ^j$. If for some $k$ and $r>0$,
\begin{equation}\label{eq:pellet}
 |p_k|r^k>\sum_{j\ne k}|p_j|r^j,
\end{equation}
then $P$ has exactly $k$ roots in $|z-a|<r$, counted with multiplicity.
\end{corollary}
\begin{proof}
On $|z-a|=r$, the difference between $P(z)$ and $p_k(z-a)^k$ has modulus
at most the right side of \eqref{eq:pellet}. Apply
\cref{thm:orderrouche} with the latter monomial as the reference polynomial.
The case $k=0$ has zero roots and is included.
\end{proof}

\begin{example}[Testing only ordinary circle points is insufficient]\label{ex:circlewarning}
Let $t=\omega^{-1}$, put
\[
 b=\frac{1+it}{\sqrt{1+t^2}},\qquad
 \alpha=(1-t^3)b,\qquad \beta=(1+t^3)b,
\]
and let $P(X)=X-\alpha$, $Q(X)=X-\beta$. Then $|b|=1$,
$|\alpha|=1-t^3<1$, and $|\beta|=1+t^3>1$. Nevertheless
$|Q(\zeta)-P(\zeta)|<|P(\zeta)|$ at every \emph{ordinary} complex unit
point $\zeta$. For $\zeta\ne1$, $v(P(\zeta))=0$; at $\zeta=1$ it is $1$,
since the leading displacement of $b$ is $it$. The error has valuation
$3$ in both cases. At the nonordinary unit point $b$, however,
$|P(b)|=t^3$ and $|Q(b)-P(b)|=2t^3$, so the required hypothesis fails.
The exact quantifier in \eqref{eq:orderrouchehyp} cannot be replaced by
sampling ordinary points.
\end{example}

\section{Exact root loci under coefficient uncertainty}\label{sec:pseudozeros}
The next result gives a sharp description, not just a sufficient error
bound. It also exposes a useful contrast between the two geometries.
Weighted coefficient-uncertainty root regions are often called
\emph{pseudozero sets} in ordinary polynomial analysis; root neighborhoods
of this type have established applications \cite{FaroukiHan}. We give the
complete argument for the surcomplex settings used here.

Fix a monic degree-$n$ polynomial $P$ and a set of perturbable indices
$J\subseteq\{0,\ldots,n-1\}$. We consider
\[
 Q(X)=P(X)+\sum_{j\in J}e_jX^j.
\]
Thus all perturbations retain the degree and leading coefficient.

\begin{theorem}[Exact uncertainty root sets in two geometries]\label{thm:pseudozeros}
Fix $z\in\SC$.

\emph{Order-modulus version.} Given $\eta_j\in\No_{\ge0}$, there exists
such a $Q$ with $Q(z)=0$ and $|e_j|\le\eta_j$ for every $j$ if and only if
\begin{equation}\label{eq:orderpseudozero}
 |P(z)|\le\sum_{j\in J}\eta_j|z^j|.
\end{equation}

\emph{Valuation version.} Given $\nu_j\in\Gamma$, there exists such a $Q$
with $Q(z)=0$ and $v(e_j)\ge\nu_j$ for every $j$ if and only if
\begin{equation}\label{eq:valpseudozero}
 v(P(z))\ge\min_{j\in J}\bigl(\nu_j+v(z^j)\bigr).
\end{equation}
Here $z^0=1$, $v(0)=+\infty$, and the minimum of an empty set is
$+\infty$. If the minimum on the right is finite, replacing every
coefficient bound by $v(e_j)>\nu_j$ replaces $\ge$ in
\eqref{eq:valpseudozero} by $>$.
\end{theorem}
\begin{proof}
The triangle inequality gives necessity of \eqref{eq:orderpseudozero}.
Let its right side be $S$. If $S=0$, the inequality forces $P(z)=0$, and
all errors may be zero. If $S>0$, put, for $z^j\ne0$,
\begin{equation}\label{eq:phaseerrors}
 e_j=-\frac{P(z)}S\,\eta_j\frac{\overline{z^j}}{|z^j|},
\end{equation}
and set $e_j=0$ for the other terms. Then $|e_j|\le\eta_j$, and
\[
 \sum_{j\in J}e_jz^j
 =-\frac{P(z)}S\sum_{j\in J}\eta_j|z^j|=-P(z).
\]
This proves sufficiency, including $z=0$ when the convention for $z^0$ is
used. The phases in \eqref{eq:phaseerrors} are algebraic unit-modulus
factors; no global argument or trigonometric branch is needed.

For the valuation version, the valuation inequality for a finite sum gives
necessity. If the minimum is finite, choose an index $j_0$ attaining it
and take
\[
 e_{j_0}=-P(z)/z^{j_0},\qquad e_j=0\quad(j\ne j_0).
\]
The assumed bound says exactly that this nonzero candidate coefficient,
when it is nonzero, has the required valuation. If the minimum is
$+\infty$, all perturbation monomials vanish at $z$ or there are none, so
$Q(z)=P(z)$; the condition correctly reduces to $P(z)=0$.
For finite minimum the same reasoning with strict inequalities proves the
last assertion. Discard monomials vanishing at $z$ when checking
necessity in that case: the minimum of finitely many quantities strictly above their respective lower bounds is
strictly above the minimum of those bounds.
\end{proof}

The order formula uses a \emph{sum} of available magnitudes; aligned
phases allow several coefficients to cooperate. The valuation formula
uses a \emph{minimum}, and one suitably chosen coefficient always suffices.
This is an exact distinction, not an analogy based on a heuristic
ultrametric norm. The proof permits infinite weights and arbitrary
surreal centers because every operation is finite.

\begin{example}[Constant-term perturbations]\label{ex:constantuncertainty}
For $P(X)=X^m-t$ with $m\ge1$, a permitted constant error of modulus at
most $\eta$ gives exactly the region $|z^m-t|\le\eta$. A constant error of
valuation at least $\nu$ gives exactly the region $v(z^m-t)\ge\nu$.
When $\nu>1$ the root residues at scale $1/m$ cannot change, although their
higher terms can. At $\nu=1$, the choice $e_0=t$ changes the polynomial to
$X^m$ and makes all roots collide at zero. The comparison anticipates the
strictness in the stability theorem of \cref{sec:stability}.
\end{example}

\section{Reduction at every center and scale}\label{sec:reduction}

\subsection{The weighted Gauss valuation}
Write a nonzero polynomial about an arbitrary surcomplex center:
\[
 P(a+Z)=\sum_{j=0}^n p_jZ^j.
\]
For $\delta\in\Gamma$ define
\begin{equation}\label{eq:gaussweight}
 w_{a,\delta}(P)=\min_{0\le j\le n}\bigl(v(p_j)+j\delta\bigr)=:w,
 \qquad
 F_{a,\delta}(U)=t^{-w}P(a+t^\delta U).
\end{equation}
We also set $w_{a,\delta}(0)=+\infty$.
At least one coefficient of $F_{a,\delta}$ has valuation zero, and all
have nonnegative valuation. Define its nonzero \emph{reduction polynomial}
by
\begin{equation}\label{eq:reductionpoly}
 R_{a,\delta}(U)=\red\bigl(F_{a,\delta}(U)\bigr)\in\C[U].
\end{equation}
It need not be monic and may have degree smaller than $P$. Terms of
strictly larger weighted valuation disappear under reduction.

\begin{lemma}[Multiplicativity and root expression]\label{lem:gauss}
For nonzero polynomials,
\begin{equation}\label{eq:gaussmult}
 w_{a,\delta}(PQ)=w_{a,\delta}(P)+w_{a,\delta}(Q).
\end{equation}
If $P(X)=a_n\prod_{i=1}^n(X-\alpha_i)$, with roots listed with
multiplicity, then
\begin{equation}\label{eq:gaussroots}
 w_{a,\delta}(P)=v(a_n)+\sum_{i=1}^n
                       \min\bigl(\delta,v(\alpha_i-a)\bigr).
\end{equation}
The convention $\min(\delta,+\infty)=\delta$ includes roots at $a$.
\end{lemma}
\begin{proof}
Normalize $P$ and $Q$ as in \eqref{eq:gaussweight}. Their reductions are
nonzero polynomials over $\C$, whose product is nonzero. Therefore the
product of the two normalized polynomials has least coefficient valuation
zero. This is exactly \eqref{eq:gaussmult}. A linear factor
$X-\alpha_i$ has weighted valuation
$\min(v(a-\alpha_i),\delta)$. Apply multiplicativity to the finite
factorization.
\end{proof}

No cancellation assumption concerning the original coefficients is
needed. Cancellation is already accounted for by computing their actual
valuations and then multiplying nonzero reduction polynomials.

\subsection{An exact ball and residue count}
\begin{theorem}[All-scale reduction counts]\label{thm:reductioncount}
Let $R=R_{a,\delta}$ be as in \eqref{eq:reductionpoly}. The following
statements count roots with multiplicity:
\begin{align}
 \#\{\alpha:P(\alpha)=0,\ v(\alpha-a)\ge\delta\}&=\deg R,
                                                        \label{eq:closedcount}\\
 \#\{\alpha:P(\alpha)=0,\ v(\alpha-a)>\delta\}&=\ord_0R.
                                                        \label{eq:opencount}
\end{align}
For every $c\in\C$, the number of roots in
$a+t^\delta(c+\mv)$ is the multiplicity of $c$ as a root of $R$.

Equivalently, if
\[
 I=\{j:v(p_j)+j\delta=w_{a,\delta}(P)\},
\]
then the closed-ball count is $\max I$, the open-ball count is $\min I$,
and the shell $v(\alpha-a)=\delta$ contains $\max I-\min I$ roots.
\end{theorem}
\begin{proof}
Put $r_i=(\alpha_i-a)/t^\delta$. In the scaled factorization, a root with
$v(r_i)<0$ contributes
\[
 U-r_i=(-r_i)(1-U/r_i).
\]
After normalizing the scalar coefficient, this factor has nonzero constant
reduction and no finite residue root. A root with $v(r_i)\ge0$ contributes
the integral monic factor $U-r_i$, whose reduction is
$U-\red(r_i)$. Multiplicativity in \cref{lem:gauss} shows that the full
normalizing scalar is a unit after its valuation is removed. Hence for
some $c_0\in\C^\times$,
\begin{equation}\label{eq:reducedfactorization}
 R(U)=c_0\prod_{v(r_i)\ge0}\bigl(U-\red(r_i)\bigr).
\end{equation}
Its degree is the number of integral $r_i$, its order at zero is the
number of infinitesimal $r_i$, and the multiplicity at any other residue
$c$ has the stated meaning. This proves the first assertions.

A coefficient of $U^j$ survives in $R$ precisely when $j\in I$.
Therefore $\deg R=\max I$ and $\ord_0R=\min I$. Subtract the two counts
for the shell formula.
\end{proof}

The theorem simultaneously allows $a$ to be infinite, $t^\delta$ to be
infinite or infinitesimal, and $\Gamma$ to have arbitrary rank. It is a
finite valued-field argument, not an appeal to analytic compactness of a
ball. A single reduction gives more information than an unlabelled
Newton polygon: it also distinguishes directions within a scale.

\begin{corollary}[Valuation Rouch\'e]\label{cor:valrouche}
If $P,Q$ are nonzero and
\begin{equation}\label{eq:valrouche}
 w_{a,\delta}(Q-P)>w_{a,\delta}(P),
\end{equation}
then their normalized reduction polynomials are equal. They have equal
root counts in the closed ball, open ball, and each residue subball
$a+t^\delta(c+\mv)$. The degrees of $P$ and $Q$ need not be equal.
\end{corollary}
\begin{proof}
After division by $t^{w_{a,\delta}(P)}$, every coefficient of the scaled
difference is infinitesimal. Thus $Q$ has the same least weight and the
same reduction as $P$. Apply \cref{thm:reductioncount}.
\end{proof}

The modulus and valuation versions of Rouch\'e have different hypotheses
and different balls. Neither is being identified with the contour theorem
on an ordinary thickening in \cite{SourceAnalysis}. The coefficient test
\eqref{eq:valrouche} has the benefit of not quantifying over boundary
points at all.

\subsection{Newton polygons over an arbitrary ordered value group}
Consider the finitely many points $(j,v(p_j))$ with $p_j\ne0$.
A lower face of slope $-\delta$ is specified algebraically by the indices
minimizing $v(p_j)+j\delta$. Slopes are ratios of value-group differences
to nonzero integers. They exist in the divisible group $\Gamma$.
This definition works even when $\Gamma$ cannot be embedded in $\R$;
we do not require a literal Euclidean drawing of its vertical axis.

\begin{corollary}[Newton polygon root valuations]\label{cor:newtonpolygon}
The finite valuations $v(\alpha_i-a)$ are exactly the negatives of the
lower-face slopes. A face whose extreme indices are $j_0<j_1$ contributes
$j_1-j_0$ roots of valuation $\delta$. A factor $(X-a)^m$ accounts
separately for $m$ roots of valuation $+\infty$.
\end{corollary}
\begin{proof}
By \cref{thm:reductioncount}, the difference of the extreme minimizing
indices at $\delta$ is exactly the number of roots with valuation
$\delta$. Thus a nontrivial face occurs exactly at a finite root valuation,
with the claimed horizontal length. Alternatively,
\eqref{eq:gaussroots} expresses the minimum of the coefficient affine
functions as a finite sum of $\min(\delta,\lambda_i)$: its slope changes
at precisely those $\lambda_i$, with the same multiplicities. Both
arguments are finite and use only the ordered group operations.
\end{proof}

\begin{example}[Two root scales and two critical directions]\label{ex:cubicnewton}
For
\begin{equation}\label{eq:cubicnewton}
 P(X)=X^3-t^2X-t^5,
\end{equation}
the nonzero coefficient points are $(0,5),(1,2),(3,0)$. Its lower faces
have slopes $-3$ and $-1$, so there is one root of valuation $3$ and two
of valuation $1$. At scale $\delta=1$,
\[
 t^{-3}P(tU)=U^3-U-t^2,\qquad R(U)=U(U-1)(U+1).
\]
The three roots are thus in the residue directions $0,1,-1$. The zero
direction has higher valuation and is resolved at the next relevant scale:
\[
 t^{-5}P(t^3V)=t^4V^3-V-1,\qquad R_{0,3}(V)=-V-1.
\]
The simple-root support recursion proved in \cref{sec:hensel} gives
\begin{align}
 \alpha_+&=t+\tfrac12t^3-\tfrac38t^5+O(t^7),\\
 \alpha_-&=-t+\tfrac12t^3+\tfrac38t^5+O(t^7),\\
 \alpha_0&=-t^3-t^7+O(t^{11}).
\end{align}
Here $O(t^k)$ means an integer-power formal Hahn series supported in
exponents at least $k$, not a fine-topological limiting estimate.
The critical points are exactly $\pm t/\sqrt3$, because
$P'=3X^2-t^2$. They are not in any of the three root submonads at scale
$1$, although both lie in their common valuation ball.
\end{example}

\begin{example}[A coefficient with nondiscrete support]\label{ex:nondiscrete}
Let
\[
 A=\sum_{k\ge0}t^{\gamma_k},\qquad
 \gamma_k=1-\frac1{k+2},\qquad P(X)=X^2-A.
\]
The support is increasing and well ordered, although it has infinitely
many exponents below $1$. Since $v(A)=1/2$, both roots have valuation
$1/4$. At that scale the reduction is $U^2-1$, so the roots have residues
$1$ and $-1$. More explicitly, write
\[
 A=t^{1/2}(1+B),\qquad
 B=\sum_{k\ge1}t^{\gamma_k-1/2}.
\]
Then the two exact roots are
\begin{equation}\label{eq:binomialhahn}
 \alpha_\pm=\pm t^{1/4}\sum_{j\ge0}\binom{1/2}{j}B^j.
\end{equation}
The binomial identity is a formal power-series identity evaluated by
\cref{lem:support}. Its use does not require that the original coefficient
series converge after replacing $t$ by a positive real number; in fact its
summands would then fail to tend to zero. Nor is the support an ordinary
integer-power series in one monomial.
\end{example}

\section{Critical-point clusters and discriminant geometry}\label{sec:critical}

\subsection{An exact non-Archimedean Rolle count}
The residue characteristic matters for differentiation. In our case the
residue field is $\C$, so every positive integer has valuation zero. This
makes the following exact count possible. Non-Archimedean Rolle and
ramification theorems have a substantial literature, including versions
for rational functions over Berkovich fields \cite{Faber}. Our proof is a
direct arbitrary-rank polynomial specialization and uses no Berkovich
space.

\begin{theorem}[Critical points in a root-containing valuation ball]\label{thm:criticalballs}
Let $P$ be nonconstant. If $\Bge{a}{\delta}$ contains $m\ge1$ roots of $P$,
counted with multiplicity, it contains exactly $m-1$ roots of $P'$.
The same assertion holds for $\Bgt{a}{\delta}$ when that open ball contains
$m\ge1$ roots. More precisely, if $R=R_{a,\delta}$ has positive degree,
the residue directions and their critical multiplicities inside the
closed ball are given by $R'$.
\end{theorem}
\begin{proof}
Let $F(U)=t^{-w}P(a+t^\delta U)$ and $R=\red(F)$. If the closed ball
contains $m\ge1$ roots, \cref{thm:reductioncount} gives $\deg R=m$.
Coefficientwise differentiation commutes with reduction, and
$R'\ne0$ because the residue characteristic is zero. Therefore $F'$ has
least coefficient valuation zero and reduction $R'$. Since
\[
 F'(U)=t^{\delta-w}P'(a+t^\delta U),
\]
the nonzero scalar has no effect on its zeros. The reduction theorem
applied to $F'$ gives a closed-ball count $\deg R'=m-1$ and the stated
residue allocation.

If the open ball contains $m\ge1$ roots, then $\ord_0R=m$. In particular
$\deg R\ge1$, so the preceding normalization still applies.
Write $R(U)=U^mS(U)$ with $S(0)\ne0$. Then
$R'(U)=U^{m-1}(mS(U)+US'(U))$, whose final factor is nonzero at zero.
Thus the open-ball count for $P'$ is $m-1$.
\end{proof}

\begin{remark}[A root-free ball is a different case]
One must not substitute $m=0$ into the theorem. The open infinitesimal
ball at zero contains no root of $X^2-1$ but contains the critical point
zero. Such critical points occur in residue directions between distinct
root branches, as the next formula makes explicit. In positive residue
characteristic even the nonempty-ball argument can fail because $R'$ may
lose its expected degree or vanish identically.
\end{remark}

\subsection{The allocation at a branching scale}
\begin{corollary}[Root branches and transverse critical directions]\label{cor:branchcritical}
Suppose the reduction in a root-containing closed ball is
\[
 R(U)=c\prod_{j=1}^s(U-c_j)^{m_j},\qquad
 c\ne0,\quad c_j\ne c_k\ (j\ne k),\quad \sum_jm_j=m.
\]
There are exactly $m_j-1$ critical points, counted with multiplicity, in
each root subball $a+t^\delta(c_j+\mv)$. The remaining $s-1$ critical
points in the parent ball have residue directions given by the roots of
\begin{equation}\label{eq:branchpoly}
 B(U)=\sum_{j=1}^s m_j\prod_{k\ne j}(U-c_k).
\end{equation}
None of these directions equals any $c_j$.
\end{corollary}
\begin{proof}
Differentiate the finite product:
\[
 R'(U)=c\prod_{j=1}^s(U-c_j)^{m_j-1}B(U).
\]
The leading coefficient of $B$ is $\sum_jm_j=m\ne0$, so its degree is
$s-1$. Also
$B(c_j)=m_j\prod_{k\ne j}(c_j-c_k)\ne0$.
Apply the directional part of \cref{thm:criticalballs}. The roots of $B$
may themselves repeat, and are counted with that multiplicity.
\end{proof}

The formula can be applied again to any smaller root-containing ball.
At every branching scale, $m_j-1$ critical points remain in each child
root cluster, while $s-1$ are accounted for at that branching. This gives
an exact finite hierarchy for a finite set of distinct roots, because
only finitely many pairwise distance valuations occur. It does not require
an infinite Newton process to separate already distinct field elements.

\subsection{Resultants and a layer formula for the discriminant}
For polynomials $P=a_n\prod_{i=1}^n(X-\alpha_i)$ and
$Q=b_m\prod_{j=1}^m(X-\beta_j)$ define their resultant by the standard
finite product
\begin{equation}\label{eq:resultantproduct}
 \Res(P,Q)=a_n^m\prod_{i=1}^nQ(\alpha_i)
 =a_n^m b_m^n\prod_{i,j}(\alpha_i-\beta_j).
\end{equation}
It is symmetric up to $(-1)^{nm}$, is a polynomial in the coefficients,
and equals the Sylvester determinant. One way to see the coefficient and
determinant assertions is first to take $P$ monic and represent
multiplication by $Q$ on $K[X]/(P)$ in the power basis. In the local
factors of \cref{thm:CRT} the determinant is
$\prod_iQ(\alpha_i)$, with roots repeated; reduction by a monic polynomial
makes each matrix entry a polynomial in the coefficients. The ordinary
Sylvester elimination identity gives the same determinant, and restoring
the leading coefficient gives \eqref{eq:resultantproduct}. We shall use
the monic multiplication-determinant definition in \cref{sec:residue}.

For a degree-$n$ polynomial define
\begin{equation}\label{eq:discdef}
 \Disc(P)=a_n^{2n-2}\prod_{i<j}(\alpha_i-\alpha_j)^2.
\end{equation}
Thus the discriminant vanishes exactly at a repeated root. For monic
$P$ one has
\[
 \Disc(P)=(-1)^{n(n-1)/2}\Res(P,P').
\]
If the relevant products do not vanish, valuation gives the exact identities
\begin{align}
 v(\Res(P,Q))&=m v(a_n)+n v(b_m)
                      +\sum_{i,j}v(\alpha_i-\beta_j),\label{eq:resval}\\
 v(\Disc(P))&=(2n-2)v(a_n)
                      +2\sum_{i<j}v(\alpha_i-\alpha_j).\label{eq:discval}
\end{align}
These quantities record collisions across all scales, not merely ordinary
coincidence of standard parts.

\begin{proposition}[Discriminant as a finite cluster sum]\label{prop:clusterdisc}
Let $P$ be monic with distinct finite roots. List the distinct positive
pairwise distance valuations as
$0<\lambda_1<\cdots<\lambda_r$ and put $\lambda_0=0$.
For each $k$, partition the roots by
\[
 \alpha_i\sim_k\alpha_j
 \quad\Longleftrightarrow\quad
 v(\alpha_i-\alpha_j)\ge\lambda_k.
\]
Let $\mathcal C_k$ be its collection of equivalence classes. Then
\begin{equation}\label{eq:clusterdisc}
 v(\Disc P)=
 2\sum_{k=1}^r(\lambda_k-\lambda_{k-1})
          \sum_{C\in\mathcal C_k}\binom{|C|}{2}.
\end{equation}
If there are no positive pairwise valuations, the sum is zero.
\end{proposition}
\begin{proof}
The ultrametric inequality makes $\sim_k$ an equivalence relation. Since
the roots are finite, all their pairwise distance valuations are
nonnegative. A pair with distance valuation $\lambda_j$ belongs to a
common class at levels $k\le j$ and not at higher levels. Its contribution
to the right side of \eqref{eq:clusterdisc}, before the factor $2$, is
\[
 \sum_{k=1}^j(\lambda_k-\lambda_{k-1})=\lambda_j.
\]
A pair of valuation zero contributes nothing. Sum this identity over the
finitely many pairs and use \eqref{eq:discval}.
\end{proof}

The formula is meaningful in any ordered divisible value group. It is a
finite telescoping identity, not an integral with respect to a real-valued
radius. It also explains why discriminant precision can protect a whole
nested root configuration.

\section{Support-controlled factor lifting}\label{sec:hensel}
The preceding sections use finite algebra after the coefficients have been
given. Constructing factors from a perturbation may require infinitely
many Hahn coefficients. Here the support convention becomes substantive.
The argument is a coprime-factor version of the preparation recursions in
\cite{SourceAnalysis,SourceGlobal}; its parameter form parallels the
fixed-domain constructions in \cite{SourceResearch}.

\subsection{The ordinary finite linear splitting}
Let
\[
 P_0=A_1\cdots A_s\in\C[X],\qquad
 A_j\text{ monic of degree }d_j,\quad \sum_jd_j=n,
\]
with the $A_j$ pairwise coprime. Consider the complex-linear map
\begin{equation}\label{eq:factorlinear}
 L:\bigoplus_{j=1}^s\C[X]_{<d_j}\longrightarrow\C[X]_{<n},
 \qquad
 (B_j)_j\longmapsto\sum_j B_j\prod_{k\ne j}A_k.
\end{equation}
It is an isomorphism. In fact, modulo $A_j$, the product of the other
$A_k$ is invertible, so the equation $L(B)=H$ uniquely specifies $B_j$
as its degree-$<d_j$ remainder. The resulting expression $L(B)-H$ is
divisible by all $A_j$, hence by $P_0$, and has degree less than $n$;
therefore it is zero. This is the precise linear operation used at each
Hahn exponent.

\subsection{Coprime Hensel factorization with a support bound}
\begin{theorem}[Support-controlled coprime factor lifting]\label{thm:hensel}
Let $P=P_0+\Delta\in\C((t^\Gamma))[X]$ be monic of degree $n$, with
$\deg\Delta<n$ and all coefficients of $\Delta$ infinitesimal. Let $E$
be the union of their positive supports and let $M=E^*$. For the pairwise
coprime monic factorization of $P_0$ above, there is a unique factorization
\begin{equation}\label{eq:henselfactors}
 P=\prod_{j=1}^s(A_j+B_j),\qquad
 \deg B_j<d_j,
\end{equation}
in which every coefficient of every $B_j$ has positive support. The
constructed supports all lie in $M\setminus\{0\}$.
Divisibility of $\Gamma$ is not required for this theorem.
\end{theorem}
\begin{proof}
Expand each unknown correction as $B_j=\sum_{\nu>0}B_{j,\nu}t^\nu$.
Proceed by well-founded recursion through $M\setminus\{0\}$.
At exponent $\nu$, the terms linear in the unknown coefficients are
$L((B_{j,\nu})_j)$. All other nonconstant terms in the product contain
at least two positive corrections. Their exponents sum to $\nu$, and
each is strictly smaller than $\nu$, so those terms are already known.
Let their sum be $C_\nu\in\C[X]_{<n}$. Define
\begin{equation}\label{eq:henselrecursion}
 (B_{j,\nu})_j=L^{-1}(\Delta_\nu-C_\nu).
\end{equation}
Each coefficient equation is finite: there are only finitely many factors,
and the support lemma leaves finitely many exponent decompositions.
All factors remain monic with their prescribed degrees. The union of the
constructed supports is contained in the well-ordered set $M$.
Equation \eqref{eq:henselrecursion} verifies the desired product at every
exponent, proving existence as an identity of Hahn polynomials.

For uniqueness, compare two such factorizations and take the least
exponent at which any correction differs, using the finite union of their
well-ordered supports. At that exponent all product terms containing
additional positive corrections use only smaller differences, hence
vanish. The remaining equation is $L$ applied to the vector of first
differences equals zero. Its invertibility contradicts the choice of that
exponent. This proves uniqueness among all positive-support candidates,
not only those preassigned to $M$.
\end{proof}

\begin{corollary}[Simple root expansion]\label{cor:simpleroot}
Every simple root $c$ of $P_0$ lifts to a unique root $c+\eta$ of $P$
with $\eta\in\mv$. Its support is contained in $E^*\setminus\{0\}$.
At the least positive error exponent $e$, its coefficient is
\begin{equation}\label{eq:firstrootcorrection}
 \eta_e=-\frac{\Delta_e(c)}{P_0'(c)}.
\end{equation}
There may be cancellation making this coefficient zero.
\end{corollary}
\begin{proof}
Apply \cref{thm:hensel} to $A_1=X-c$ and $A_2=P_0/(X-c)$, which are
coprime. The lifted first factor is $X-c-\eta$. For any root reducing to
$c$, the reduction theorem, or the uniqueness of this linear factor,
identifies it with $c+\eta$. Expanding $P(c+\eta)=0$ at the least positive
input exponent gives
$P_0'(c)\eta_e+\Delta_e(c)=0$.
\end{proof}

The coprimality assumption is essential for this support conclusion.
For $P=X^m-t$, the reduction $X^m$ has a repeated root. Its individual
roots require $t^{1/m}$, which need not lie in the original exponent
group or in the monoid generated by $1$. The theorem nevertheless lifts
any decomposition into \emph{distinct coprime clusters} in the original
coefficient field. Symmetric cluster coefficients can have substantially
simpler supports than their separately labelled roots.

\subsection{Holomorphic parameters on a fixed domain}
\begin{theorem}[Parameter factor lifting without repeated shrinking]\label{thm:parameterhensel}
Let $U\subseteq\C^d$ be an ordinary domain. Suppose
$A_1(x,X),\ldots,A_s(x,X)\in\OO(U)[X]$ are monic of fixed degrees and are
pairwise coprime at every $x\in U$. Let
\[
 P=\prod_jA_j+\Delta\in\mathcal H_{\Gamma,\ge0}(U)[X]
\]
be monic of the same total degree, with $\Delta$ of strictly positive
support and lower degree. If $E$ is its total positive coefficient
support, the factors in \eqref{eq:henselfactors} exist uniquely with
$B_j\in\mathcal H_{\Gamma,>0}(U)[X]$, $\deg B_j<d_j$, and supports
contained in $E^*\setminus\{0\}$. Every ordinary coefficient function is
holomorphic on the \emph{original} domain $U$.
\end{theorem}
\begin{proof}
In the finite power bases, \eqref{eq:factorlinear} is a square matrix
$L(x)$ of ordinary holomorphic functions. Its value is invertible at every
$x$ by the preceding coprime argument. Consequently
$L(x)^{-1}=\operatorname{adj}(L(x))/\det L(x)$ is holomorphic throughout
$U$. There is no requirement that its ordinary size be uniformly bounded.
Perform \eqref{eq:henselrecursion} with this fixed holomorphic inverse.
Finite products and applications of $L^{-1}$ preserve holomorphy on the
same $U$ and introduce no Hahn exponent. The support and uniqueness
arguments are unchanged.
\end{proof}

When the leading ordinary polynomial has only simple roots, the same
argument on its ordinary root cover gives a unique coherent lift of every
sheet. Local lifts agree on overlaps by uniqueness. Ordinary analytic
continuation therefore permutes the lifted sheets exactly as it permutes
the original ones, as also proved in \cite{SourceResearch}. This does not
produce a globally labelled root on a base with nontrivial ordinary
monodromy, and does not assert holomorphic individual roots through a
collision. The factor theorem keeps precisely the hypotheses that make
its finite linear splitting invertible.

\section{Quantitative preservation of the root configuration}\label{sec:stability}

\subsection{A root-specific precision threshold}
Let $P$ be monic of degree $n\ge2$, with distinct finite roots
$\alpha_1,\ldots,\alpha_n$. Write
\begin{equation}\label{eq:stabilitydata}
 d_{ij}=v(\alpha_i-\alpha_j),\qquad
 \sigma_i=v(P'(\alpha_i))=\sum_{j\ne i}d_{ij},\qquad
 \delta_i=\max_{j\ne i}d_{ij}.
\end{equation}
All these values are nonnegative. The number $\delta_i$ records the
closest other root in valuation geometry; larger valuation means smaller
distance. The open balls $\Bgt{\alpha_i}{\delta_i}$ are pairwise disjoint.
Indeed, a point in two of them would force
$v(\alpha_i-\alpha_j)>\min(\delta_i,\delta_j)\ge d_{ij}$.

\begin{theorem}[Exact matching and full cluster stability]\label{thm:stability}
Let
\[
 Q=P+H,\qquad H(X)=\sum_{j=0}^{n-1}h_jX^j,
 \qquad \mu=\min_jv(h_j),
\]
where $Q$ is monic of degree $n$. Suppose $H\ne0$ and
\begin{equation}\label{eq:stabilitythreshold}
 \mu>T(P):=\max_i(\sigma_i+\delta_i).
\end{equation}
Then $Q$ has $n$ distinct finite roots with a unique labelling
$\beta_1,\ldots,\beta_n$ such that
\begin{equation}\label{eq:matchingballs}
 v(\beta_i-\alpha_i)>\delta_i\qquad(1\le i\le n).
\end{equation}
Furthermore,
\begin{align}
 v(\beta_i-\beta_j)&=d_{ij} &&(i\ne j),\label{eq:pairwisepreserved}\\
 v(\beta_i-\alpha_i)&\ge\mu-\sigma_i,\label{eq:shiftbound}\\
 v(Q'(\beta_i))&=\sigma_i,\qquad
 v(\Disc Q)=v(\Disc P).\label{eq:discpreserved}
\end{align}
If $H(\alpha_i)\ne0$, the displacement has the sharper leading-term
identity
\begin{equation}\label{eq:leadingdisplacement}
 \frac{\beta_i-\alpha_i}{-H(\alpha_i)/P'(\alpha_i)}\in1+\mv,
 \qquad
 v(\beta_i-\alpha_i)=v(H(\alpha_i))-\sigma_i.
\end{equation}
If $H(\alpha_i)=0$, then $\beta_i=\alpha_i$. For $H=0$ all conclusions
hold trivially, with the zero-displacement case understood separately.
\end{theorem}
\begin{proof}
Fix $i$ and use center $\alpha_i$ and scale $\delta_i$.
Formula \eqref{eq:gaussroots} gives
\begin{equation}\label{eq:rootweight}
 w_{\alpha_i,\delta_i}(P)=\delta_i+\sum_{j\ne i}d_{ij}
 =\delta_i+\sigma_i.
\end{equation}
The open ball at this center and scale contains just the simple root
$\alpha_i$: for every other root, $d_{ij}\le\delta_i$.

Every coefficient in the translated expansion of $H(\alpha_i+Z)$ has
valuation at least $\mu$. To see this, use the finite binomial formula;
$\alpha_i$ is integral, the original errors have valuation at least
$\mu$, and integer binomial coefficients have nonnegative valuation.
Because $\delta_i\ge0$, scaling by $t^{\delta_i}$ cannot lower these
coefficient valuations. Therefore
\[
 w_{\alpha_i,\delta_i}(H)\ge\mu>
             \sigma_i+\delta_i=w_{\alpha_i,\delta_i}(P).
\]
Valuation Rouch\'e gives exactly one root of $Q$, counted with multiplicity,
in $\Bgt{\alpha_i}{\delta_i}$. The balls are disjoint, so these are $n$
distinct simple roots and exhaust the degree of $Q$. Their finiteness is
immediate from that of the centers and the nonnegative scales.
This proves the matching and its uniqueness.

Set $u_i=\beta_i-\alpha_i$. For $i\ne j$, both corrections $u_i,u_j$
have valuation strictly greater than $d_{ij}$. The term
$\alpha_i-\alpha_j$ therefore determines the leading term of
$\beta_i-\beta_j$, proving \eqref{eq:pairwisepreserved}. The derivative
and discriminant identities follow by taking valuations of their finite
root products.

For the displacement estimate, all factors
$\beta_i-\alpha_j$ with $j\ne i$ have valuation $d_{ij}$. Hence
\[
 v(P(\beta_i))=v(u_i)+\sigma_i
\]
when $u_i\ne0$. Since $P(\beta_i)=-H(\beta_i)$ and $\beta_i$ is integral,
$v(H(\beta_i))\ge\mu$, giving \eqref{eq:shiftbound}. A zero displacement
satisfies that inequality under the $+\infty$ convention.

For the leading-term assertion, if $H(\alpha_i)=0$ then $\alpha_i$ itself
is a root of $Q$ in its assigned ball, so uniqueness gives $u_i=0$.
Otherwise $u_i\ne0$. The exact factorization gives
\begin{equation}\label{eq:linearleading}
 P(\alpha_i+u_i)=P'(\alpha_i)u_i
             \prod_{j\ne i}\left(1+\frac{u_i}{\alpha_i-\alpha_j}\right)
 =P'(\alpha_i)u_i(1+\varepsilon_i),
 \quad\varepsilon_i\in\mv.
\end{equation}
Also $H(\alpha_i+u_i)-H(\alpha_i)=u_i V_i$ with $v(V_i)\ge\mu$, by the
finite difference-of-powers identity and integrality of both arguments.
The value of this difference is strictly larger than
$\sigma_i+v(u_i)$, because $\mu>\sigma_i$. It therefore cannot affect
the leading term in
$P(\alpha_i+u_i)+H(\alpha_i+u_i)=0$. Equation
\eqref{eq:linearleading} proves \eqref{eq:leadingdisplacement}, including
the exact valuation of the shift.
\end{proof}

The theorem preserves the whole matrix of pairwise valuations, so every
root partition defined by a valuation threshold is preserved. This is a
precise meaning of preservation of the complete finite cluster geometry.
It does not assert equality of the roots, of their ordinary moduli, or of
every Hahn coefficient of the pairwise differences.

\subsection{A discriminant-only sufficient condition}
\begin{corollary}[Discriminant precision protects all clusters]\label{cor:discprecision}
Under the hypotheses on $P$, put $D=v(\Disc P)$. Then $T(P)\le D$.
Consequently a monic same-degree perturbation with
\begin{equation}\label{eq:discriminantprecision}
 \min_jv(h_j)>D
\end{equation}
has all the matching and preservation properties of
\cref{thm:stability}.
\end{corollary}
\begin{proof}
By \eqref{eq:discval}, $D=2\sum_{i<j}d_{ij}$. For fixed $i$, choose
$j_0$ with $d_{ij_0}=\delta_i$. In the sum $\sigma_i+\delta_i$, the
pair $(i,j_0)$ is counted twice, every other pair incident to $i$ once,
and the remaining pairs not at all. All $d_{ij}$ are nonnegative, so this
is at most $D$. Take the maximum over $i$ and apply
\cref{thm:stability}.
\end{proof}

\begin{example}[Why the inequality must be strict]\label{ex:sharpdisc}
Let $h>0$ be a value-group element, and consider
\[
 P(X)=X(X-t^h),\qquad
 Q(X)=(X-\tfrac12t^h)^2=P(X)+\tfrac14t^{2h}.
\]
The two roots of $P$ are distinct, and
$D=T(P)=2h$. The error has valuation exactly $2h$, but $Q$ has a double
root. Thus replacing $>$ by $\ge$ in the universal criterion
\eqref{eq:discriminantprecision} is false. The example does not claim that
$D$ is the optimal threshold for every polynomial; \cref{ex:quartic}
shows that $T(P)$ can be strictly smaller.
\end{example}

\subsection{How to use the theorem at infinite scales}\label{subsec:rescaledprecision}
Finiteness of the roots in \cref{thm:stability} is a normalization
hypothesis, not a restriction to an Archimedean problem. Let $P$ be any
monic squarefree degree-$n$ polynomial. Choose $a\in\SC$ and a positive
surreal $\rho$ large enough that all $(\alpha_i-a)/\rho$ are finite.
The translated root bound in \cref{prop:rootbounds} supplies such a $\rho$
from the coefficients. Put
\[
 P_*(Y)=\rho^{-n}P(a+\rho Y),\qquad
 Q_*(Y)=\rho^{-n}Q(a+\rho Y).
\]
Both are monic of degree $n$. If $H(a+Z)=\sum_{j<n}\widetilde h_jZ^j$,
the normalized coefficient precision is
\begin{equation}\label{eq:rescaledmu}
 \mu_* =\min_{j<n}\bigl(v(\widetilde h_j)+(j-n)v(\rho)\bigr).
\end{equation}
The discriminant transforms exactly as
\begin{equation}\label{eq:rescaleddisc}
 \Disc(P_*)=\rho^{-n(n-1)}\Disc(P),\qquad
 D_* =v(\Disc P)-n(n-1)v(\rho).
\end{equation}
Thus $\mu_*>D_*$ is a coefficient-level sufficient condition for the
normalized matching. Pulling the roots back by $a+\rho Y$ preserves all
pairwise valuation equalities in the original coordinates as well.
Applying the unnormalized discriminant criterion without this adjustment
would be unjustified when some original roots are infinite.

\subsection{A complete four-root hierarchy}\label{subsec:quartic}
\begin{example}[A stronger threshold than the discriminant alone]\label{ex:quartic}
Let
\begin{equation}\label{eq:quartic}
 P(X)=(X-t)(X-t-t^3)(X+t)(X-1),\qquad t=\omega^{-1},
\end{equation}
and label its roots in that order. The positive distance valuations are
\[
 d_{12}=3,\qquad d_{13}=d_{23}=1;
 \qquad d_{14}=d_{24}=d_{34}=0.
\]
Hence
\begin{equation}\label{eq:quarticdata}
 (\sigma_1,\sigma_2,\sigma_3,\sigma_4)=(4,4,2,0),\qquad
 (\delta_1,\delta_2,\delta_3,\delta_4)=(3,3,1,0),
\end{equation}
and $T(P)=7$, whereas $D=10$. The discriminant has leading term
$16t^{10}$. Formula \eqref{eq:clusterdisc} gives the same valuation as
\[
 2\left(1\binom32+(3-1)\binom22\right)=10.
\]

The critical points can be placed without solving a cubic. At scale zero
about zero, the reduction is $U^3(U-1)$, whose derivative is
$U^2(4U-3)$. Two critical points therefore lie in the infinitesimal monad
of zero, and one has standard part $3/4$. At scale one about zero the
reduction, up to a nonzero scalar, is $(U-1)^2(U+1)$. Its derivative is
$(U-1)(3U+1)$, so one critical point satisfies $X/t\in1+\mv$ and another
satisfies $X/t\in-1/3+\mv$. In the smaller ball of scale three about $t$,
the reduction is a nonzero scalar times $U(U-1)$. Thus the first of these
has
\[
 (X-t)/t^3\in\tfrac12+\mv.
\]
These three counts exhaust $\deg P'=3$ and describe the successive
branching directions of every critical point.

Now set $Q=P+t^8$. Its coefficient precision $8$ exceeds $T(P)=7$ but
not $D=10$. The stronger theorem applies and preserves every pairwise
root-distance valuation. Its uniquely matched roots satisfy
\begin{align}
 \beta_1-\alpha_1&=-\tfrac12t^4(1+\varepsilon_1),&
 \beta_2-\alpha_2&= \tfrac12t^4(1+\varepsilon_2),\\
 \beta_3-\alpha_3&= \tfrac14t^6(1+\varepsilon_3),&
 \beta_4-\alpha_4&=-t^8(1+\varepsilon_4),
\end{align}
with all $\varepsilon_i\in\mv$. These leading terms follow directly from
$-t^8/P'(\alpha_i)$. They are checked independently by exact symbolic
expansion in the accompanying verification script.
\end{example}

\section{Universal residue duality and collision-stable traces}\label{sec:residue}
This section deliberately works over a more general coefficient ring.
Let $R$ be any commutative ring with identity, and let
\[
 P(X)=X^n+p_{n-1}X^{n-1}+\cdots+p_0\in R[X],\qquad n\ge1.
\]
Monic division works over $R$ without dividing by any coefficient.
Thus $A=R[X]/(P)$ is free with basis $1,X,\ldots,X^{n-1}$.
Define
\begin{equation}\label{eq:lambdadef}
 \lambda_P:A\longrightarrow R,\qquad
 \lambda_P(h)=[X^{n-1}]\remp_P h.
\end{equation}
For $h\in A$, let $m_h$ denote its multiplication endomorphism. This
algebraic functional will recover the residue and trace formulas of the
supplied manuscripts without choosing a contour or labelling moving roots.

\subsection{A dual basis with no discriminant denominator}
\begin{theorem}[Universal residue pairing and explicit dual basis]\label{thm:residuepairing}
The bilinear form $(f,g)\mapsto\lambda_P(fg)$ on $A$ is perfect.
Writing $p_n=1$, the residue-dual basis to $1,X,\ldots,X^{n-1}$ is
\begin{equation}\label{eq:dualbasis}
 b^i(X)=\sum_{k=i+1}^n p_kX^{k-1-i},\qquad 0\le i<n.
\end{equation}
In particular $\lambda_P(X^jb^i)=\delta_{ij}$. Its Gram matrix
$G_{ij}=\lambda_P(X^{i+j})$ has
\begin{equation}\label{eq:gramdet}
 \det G=(-1)^{n(n-1)/2}.
\end{equation}
The coefficient matrix of
\begin{equation}\label{eq:bezoutkernel}
 \mathcal B_P(X,Y)=\frac{P(X)-P(Y)}{X-Y}
                 =\sum_{i=0}^{n-1}b^i(X)Y^i
\end{equation}
is $G^{-1}$ in these power bases.
\end{theorem}
\begin{proof}
There is a formal Laurent expansion at infinity
\[
 \frac1{P(X)}=X^{-n}
    \bigl(1+p_{n-1}X^{-1}+\cdots+p_0X^{-n}\bigr)^{-1}.
\]
The inverse is a formal power series in the independent symbol $X^{-1}$,
valid over any $R$. For every polynomial $h$, division by $P$ gives
\begin{equation}\label{eq:lambdaatinfinity}
 \lambda_P(h)=[X^{-1}]\frac{h(X)}{P(X)}.
\end{equation}
Indeed, the quotient is polynomial and has no $X^{-1}$ term; for a
remainder of degree less than $n$, that coefficient is exactly its
$X^{n-1}$ coefficient.

The definition of $b^i$ yields
\[
 \frac{b^i(X)}{P(X)}
 =X^{-i-1}-\frac{\sum_{k=0}^i p_kX^{k-i-1}}{P(X)}.
\]
The second term has highest possible Laurent degree $-n-1$. After
multiplication by $X^j$ with $0\le j<n$, it cannot contribute to degree
$-1$. Formula \eqref{eq:lambdaatinfinity} therefore gives
$\lambda_P(X^jb^i)=\delta_{ij}$. This constructs an inverse to the map
$A\to\operatorname{Hom}_R(A,R)$ defined by the pairing, proving perfection
over the ring itself.

For the determinant, $G_{ij}=0$ when $i+j<n-1$ and $G_{ij}=1$ when
$i+j=n-1$. Reverse the order of its columns. The resulting matrix is
triangular with all diagonal entries $1$, proving \eqref{eq:gramdet}.
Finally the finite identity
$(X^k-Y^k)/(X-Y)=\sum_{i=0}^{k-1}X^{k-1-i}Y^i$ proves
\eqref{eq:bezoutkernel}. Its $i$th coefficient column is precisely the
dual vector $b^i$, so the coefficient matrix is $G^{-1}$.
\end{proof}

No inverse root separation, resultant, or discriminant appears. This is
stronger than a statement that the form is nondegenerate over the fraction
field: it is an isomorphism over the given ring $R$, even when that ring
has zero divisors or is non-Noetherian.

\subsection{Trace, norm, and the derivative as the trace element}
\begin{theorem}[Universal trace identity]\label{thm:trace}
For every $h\in A$,
\begin{equation}\label{eq:traceidentity}
 \Tr(m_h)=\lambda_P(P'h).
\end{equation}
If $G_{\mathrm{tr},ij}=\Tr(m_{X^{i+j}})$ is the trace-pairing matrix,
then
\begin{equation}\label{eq:tracegram}
 G_{\mathrm{tr}}=G M_{P'},\qquad
 \det G_{\mathrm{tr}}=\Disc(P),
\end{equation}
where $M_{P'}$ is the multiplication matrix of $P'$ and the discriminant
is the usual universal coefficient polynomial for monic $P$.
For any polynomial $Q$, $\det(m_Q)=\Res(P,Q)$.
\end{theorem}
\begin{proof}
The coefficient of $X^i$ in the power-basis expansion of $hX^i$ is
$\lambda_P(hX^ib^i)$. Summing the diagonal coefficients gives
\[
 \Tr(m_h)=\lambda_P\left(h\sum_{i=0}^{n-1}X^ib^i\right).
\]
But the explicit dual basis satisfies the polynomial identity
\begin{equation}\label{eq:eulerderivative}
 \sum_{i=0}^{n-1}X^ib^i
 =\sum_{k=1}^n\sum_{i=0}^{k-1}p_kX^{k-1}
 =\sum_{k=1}^n kp_kX^{k-1}=P'.
\end{equation}
This proves \eqref{eq:traceidentity} over every $R$.

The trace-pairing entry is $\lambda_P(X^iP'X^j)$, which is the $(i,j)$
entry of $GM_{P'}$. The determinant of multiplication by $Q$ modulo a
monic $P$ is the monic resultant. For clarity, this equality can be seen
without assuming roots over $R$: both are polynomials with integer
coefficients in the finitely many coefficients of $P,Q$. Over the
fraction field of the universal integer coefficient ring, $P$ has
distinct roots and the multiplication determinant is their evaluation
product, which is \eqref{eq:resultantproduct}. The universal coefficient
ring injects into that field, so the polynomial identity holds there and
hence after every specialization to $R$.

Taking determinants in $G_{\mathrm{tr}}=GM_{P'}$ now gives
$(-1)^{n(n-1)/2}\Res(P,P')$, equal to $\Disc(P)$ for monic $P$.
The same universal-polynomial argument identifies it with
\eqref{eq:discdef} over splitting fields and with its coefficient
specialization over arbitrary rings.
\end{proof}

The proof supplies an explicit one-variable answer to the trace-element
question: the element representing the trace under residue duality is
exactly $P'$. This should not be mistaken for a proof of the general
multivariable Jacobian comparison left separate in
\cite{SourceResearch}; our argument concerns the monogenic algebra
$R[X]/(P)$.

\begin{proposition}[Field-level traces and norms, including collisions]\label{prop:tracenormroots}
Over an algebraically closed characteristic-zero field, write
$P=\prod_i(X-\alpha_i)^{m_i}$. Then
\begin{equation}\label{eq:tracenormroots}
 \Tr(m_h)=\sum_i m_i h(\alpha_i),\qquad
 \Nm_A(h):=\det(m_h)=\prod_i h(\alpha_i)^{m_i}.
\end{equation}
The radical of the trace pairing is the nilradical of $A$. The residue
pairing of \cref{thm:residuepairing} has zero radical, whether or not roots
repeat.
\end{proposition}
\begin{proof}
Use the CRT factors $K[U]/(U^{m_i})$. On the $i$th factor, multiplication
by $h(\alpha_i+U)$ is triangular in the powers of $U$, with all diagonal
entries $h(\alpha_i)$. Its trace and determinant are the corresponding
expressions in \eqref{eq:tracenormroots}. On a pair of elements the trace
pairing is thus $\sum_i m_i f(\alpha_i)g(\alpha_i)$.
If all $f(\alpha_i)=0$, this is zero for every $g$. Conversely, choose the
CRT idempotent of the $i$th factor to deduce $m_i f(\alpha_i)=0$ from
membership in the radical. Characteristic zero gives $f(\alpha_i)=0$.
These are precisely the elements whose image in each local factor lies
in its nilpotent maximal ideal, hence precisely the nilradical. The
residue pairing is perfect by its already constructed dual basis.
\end{proof}

\subsection{Finite local residues and Hermite information}
\begin{proposition}[Residue formula at multiple roots]\label{prop:localresidues}
Let $P$ be monic over a characteristic-zero algebraically closed field,
and set $P_i=P/(X-\alpha_i)^{m_i}$. Then for every polynomial $h$,
\begin{equation}\label{eq:localresidues}
 \lambda_P(h)=\sum_i\frac1{(m_i-1)!}
 \left(\frac{d^{m_i-1}}{dX^{m_i-1}}\frac{h(X)}{P_i(X)}
 \right)_{X=\alpha_i}.
\end{equation}
In the squarefree case this becomes
\begin{equation}\label{eq:simpleresidues}
 \lambda_P(h)=\sum_{i=1}^n\frac{h(\alpha_i)}{P'(\alpha_i)}.
\end{equation}
Both formulas are finite algebraic identities for surcomplex inputs.
\end{proposition}
\begin{proof}
Finite partial fractions express $h/P$ as a polynomial plus a sum of
terms $c_{i,k}(X-\alpha_i)^{-k}$. In the Laurent expansion at infinity,
only the terms with $k=1$ contribute to $X^{-1}$, and their contributions
are $c_{i,1}$. Near $\alpha_i$, put $U=X-\alpha_i$. The coefficient
$c_{i,1}$ is the coefficient of $U^{m_i-1}$ in the Taylor expansion of
$h(\alpha_i+U)/P_i(\alpha_i+U)$. Its denominator has nonzero constant
term, so that coefficient is exactly the derivative in
\eqref{eq:localresidues}. Apply \eqref{eq:lambdaatinfinity} and sum.
For a simple root, $P_i(\alpha_i)=P'(\alpha_i)$.
\end{proof}

The local terms in \eqref{eq:simpleresidues} may individually be infinite
even when $\lambda_P(h)$ is finite or zero. This is not divergence: it is
cancellation in a finite sum. At a collision, separate point weights can
become singular while the universal functional remains defined, and its
local formula retains derivative information.

\begin{example}[A cubic collision revisited from the source]\label{ex:cubicduality}
The polynomial $P=X^3-aX-b$ used in \cite{SourceResearch} provides a useful
normalization check. In the basis $(1,X,X^2)$,
\[
 G=\begin{pmatrix}0&0&1\\0&1&0\\1&0&a\end{pmatrix},
 \qquad
 G^{-1}=\begin{pmatrix}-a&0&1\\0&1&0\\1&0&0\end{pmatrix},
 \qquad \det G=-1.
\]
The dual basis is $(X^2-a,X,1)$, whose trace element is $3X^2-a=P'$.
The trace matrix is
\[
 G_{\mathrm{tr}}=
 \begin{pmatrix}3&0&2a\\0&2a&3b\\2a&3b&2a^2\end{pmatrix},
 \qquad \det G_{\mathrm{tr}}=4a^3-27b^2.
\]
At $a=3t^2$, $b=2t^3$, one has $P=(X-2t)(X+t)^2$ and
\begin{equation}\label{eq:cubiccollision}
 \lambda_P(h)=\frac{h(2t)-h(-t)}{9t^2}-\frac{h'(-t)}{3t},
 \qquad \Tr(m_h)=h(2t)+2h(-t).
\end{equation}
The trace forgets the nilpotent direction at the double point; the residue
functional remembers its first derivative. This example is inherited and
explicitly credited, while \cref{thm:residuepairing,thm:trace} give the
universal finite-ring proof and the explicit dual basis in every degree.
\end{example}

\subsection{Coherent coefficient families}
Take $R=\HH(U)$ or $R=\mathcal H_{\Gamma,\ge0}(U)$, for a fixed ordinary
domain $U$. Every monic bounded-degree family $P(x,X)$ gives a free
rank-$n$ algebra over $R$. Polynomial division, multiplication matrices,
$\lambda_P$, the dual basis, traces, norms, and discriminants are computed
by finite coefficient-ring operations. Consequently they are legitimate
Hahn-coherent coefficient functions on the same $U$. No inverse
discriminant or repeated domain shrink is needed for the residue pairing.

For an integral family, all entries in its residue Gram matrix and in
its explicit inverse are integral. If the input coefficient supports lie
in a nonnegative support monoid $M$, these entries also lie in $M$.
At a finite surcomplex parameter, evaluation \eqref{eq:coherentevaluation}
is a ring homomorphism, so it commutes with these finite identities. The
resulting specialized pairing is exactly the pairing for the specialized
polynomial. This explains collision stability algebraically, independently
of a simultaneous choice of root labels.

For more general coherent analytic numerators of unbounded Taylor degree,
one must use the fixed-domain division theorems of the supplied manuscripts
before taking a remainder. The purely polynomial argument in this section
does not silently replace that extra analytic step.

\section{Finite polynomial maps, branch values, and ramification}\label{sec:maps}

\subsection{Every fiber has the same algebraic length}
\begin{theorem}[A polynomial is a finite flat map of its degree]\label{thm:finitemap}
Let $P\in K[X]$ be monic of degree $n\ge1$. Regard $K[X]$ as a $K[Y]$
algebra by $Y\mapsto P(X)$. It is free of rank $n$, with basis
$1,X,\ldots,X^{n-1}$. For every $y\in K$, the fiber algebra
\[
 K[X]/(P(X)-y)
\]
has dimension $n$, and decomposes as
\begin{equation}\label{eq:fiberdecomp}
 \prod_{P(\alpha)=y}K[U]/(U^{e_\alpha}),
 \qquad e_\alpha=\ord_\alpha(P-y),\qquad
 \sum_{P(\alpha)=y}e_\alpha=n.
\end{equation}
\end{theorem}
\begin{proof}
The presentation $K[Y,X]/(P(X)-Y)$ is isomorphic to $K[X]$: eliminate
$Y$ by substitution $Y=P(X)$. The defining polynomial is monic in $X$,
so division makes this algebra free over $K[Y]$ with the stated basis.
Specializing $Y$ to $y$ leaves the same basis and rank. The factorization
and CRT theorem prove \eqref{eq:fiberdecomp}.
\end{proof}

The words ``finite flat'' describe this free algebra and the associated
map of set-sized spectra. They do not assert topological properness or
path lifting in the full fine surreal topology. In particular, the
ordinary connected real-parameter paths excluded in
\cite{SourceAnalysis} play no role.

\subsection{Critical values and the discriminant polynomial}
\begin{theorem}[Branch-value polynomial]\label{thm:branchvalues}
Suppose $n\ge2$. List the roots of $P'$ as $c_1,\ldots,c_{n-1}$,
including multiplicity. Then
\begin{equation}\label{eq:branchvalues}
 \Disc_X(P(X)-Y)
 =(-1)^{n(n-1)/2}n^n\prod_{j=1}^{n-1}(P(c_j)-Y).
\end{equation}
It has degree $n-1$ in $Y$. Its zeros are exactly the values with a
nonreduced fiber, equivalently the critical values of $P$, with the
multiplicities indicated by this product.
\end{theorem}
\begin{proof}
Put $A(X)=P(X)-Y$. The monic discriminant identity gives
\[
 \Disc_X A=(-1)^{n(n-1)/2}\Res_X(A,P').
\]
The resultant symmetry sign is $(-1)^{n(n-1)}=1$, and $P'$ has leading
coefficient $n$. Hence the product definition gives
\[
 \Res_X(A,P')=\Res_X(P',A)=n^n\prod_{j=1}^{n-1}A(c_j),
\]
which is \eqref{eq:branchvalues}. A fiber has a multiple root exactly when
$A$ and its derivative have a common root, so the zero set is as claimed.
The leading coefficient in $Y$ is nonzero in characteristic zero, giving
degree $n-1$.
\end{proof}

Coincident critical values combine their multiplicities in this product;
there need not be $n-1$ distinct branch values. The formula continues to
make sense over coherent coefficient rings as a universal polynomial
identity, even when those values cannot be labelled globally.

\subsection{The polynomial ramification sum}
\begin{proposition}[Finite and infinite ramification]\label{prop:ramification}
For a nonconstant degree-$n$ polynomial in characteristic zero, the local index at a
finite point $\alpha$ is $e_\alpha=\ord_\alpha(P-P(\alpha))$, and
\begin{equation}\label{eq:finiteRH}
 \sum_{\alpha\in K}(e_\alpha-1)=n-1.
\end{equation}
Only finitely many summands are nonzero. At infinity the local index is
$n$. Thus the total sum on the algebraic projective line, including
infinity, is $2n-2$.
\end{proposition}
\begin{proof}
The finite Taylor expansion at $\alpha$ begins with a nonzero term of
degree $e_\alpha$. Differentiation therefore makes the multiplicity of
$\alpha$ in $P'$ equal to $e_\alpha-1$. Summing the multiplicities of
the degree-$n-1$ polynomial $P'$ proves \eqref{eq:finiteRH}.

For the point at infinity use the formal source coordinate $u=1/X$ and
target coordinate $v=1/Y$. If $a_n$ is the leading coefficient, then
\[
 \frac1{P(1/u)}=a_n^{-1}u^n(1+\text{higher powers of }u)^{-1}.
\]
The unit factor has nonzero constant term, so the order is $n$. Add its
contribution $n-1$ to the finite sum. This is the polynomial instance of
the ramification formula, derived directly rather than from a
topological covering or surface argument.
\end{proof}

\begin{corollary}[Surjectivity and polynomial injections]\label{cor:polynomialmaps}
Every nonconstant polynomial $\SC\to\SC$ is surjective. The injective
polynomials are exactly the affine polynomials with nonzero slope.
\end{corollary}
\begin{proof}
For every target value, algebraic closure supplies a root of $P-y$.
If $n\ge2$, choose $y$ outside the finite set of critical values supplied
by \cref{thm:branchvalues}. Its fiber has $n$ distinct roots, so the map is
not injective. The affine case is immediate. This proof is finite and
does not require an analytic inverse theorem.
\end{proof}

\section{Several-variable polynomial equations and class-safe certificates}\label{sec:nullstellensatz}
The field-based algebra in several variables remains conventional, but
its scope must be stated carefully. We prove the finite-system form of
the Nullstellensatz and explain what does not follow about arbitrary
Hahn-coherent analytic algebras. The standard underlying theorem is
recorded, for example, in the Stacks Project \cite[Tag~00FV]{StacksNS}.

\subsection{The finite-type field lemma}
\begin{lemma}[Zariski's lemma in the needed form]\label{lem:zariski}
If an extension field $L$ of an infinite field $k$ is finitely generated
as a $k$-algebra, then it is finite algebraic over $k$.
\end{lemma}
\begin{proof}
Choose a maximal algebraically independent subcollection
$t_1,\ldots,t_r$ among a finite algebra generating set. The other
finitely many generators are algebraic over $k(t_1,\ldots,t_r)$, so
$L$ is a finite field extension of that rational-function field.
Clear the coefficients' denominators in their monic algebraic equations:
for some nonzero $f\in k[t_1,\ldots,t_r]$, every remaining generator is
integral over
\[
 A=k[t_1,\ldots,t_r,1/f].
\]
Because $L$ was generated as an algebra by the original generators and is
a field containing $1/f$, it equals the algebra obtained by adjoining
those integral generators to $A$. Thus $L$ is integral over $A$.

If a field is integral over a subring $A$, that subring is a field. Indeed,
for $0\ne a\in A$, the element $a^{-1}\in L$ satisfies a monic equation;
multiplying the equation by a suitable power of $a$ expresses $a^{-1}$
as an element of $A$.

If $r>0$, choose $c\in k$ such that $t_1-c$ does not divide $f$; only
finitely many such linear factors can divide a fixed nonzero polynomial.
Then $1/(t_1-c)\notin A$, as can be checked by multiplying a putative
representation by a power of $f$ and using divisibility in
$k[t_1,\ldots,t_r]$. This contradicts that $A$ is a field. Hence $r=0$,
and the finitely many original generators are algebraic, giving a finite
extension.
\end{proof}

\subsection{Finite surcomplex systems}
\begin{theorem}[Surcomplex Nullstellensatz for finite data]\label{thm:nullstellensatz}
Let $F_1,\ldots,F_s,G\in\SC[X_1,\ldots,X_d]$ be finitely many
finite-degree polynomials. Then $G$ vanishes at every common surcomplex
zero of the $F_i$ if and only if, for some ordinary integer $N\ge1$ and
finite polynomials $H_i\in\SC[X_1,\ldots,X_d]$,
\begin{equation}\label{eq:nullcertificate}
 G^N=\sum_{i=1}^sH_iF_i.
\end{equation}
In particular the $F_i$ have no common surcomplex zero if and only if
$1$ belongs to their polynomial ideal. The certificate can be chosen in
any algebraically closed Hahn workspace containing all coefficients.
\end{theorem}
\begin{proof}
Fix such a workspace $K$. A maximal ideal $\mathfrak n$ of the set ring
$K[X_1,\ldots,X_d]$ gives a field quotient that is finitely generated as
a $K$-algebra. By \cref{lem:zariski} and algebraic closure it is $K$.
The images of the variables are coordinates $a_1,\ldots,a_d\in K$,
and $\mathfrak n=(X_1-a_1,\ldots,X_d-a_d)$. Consequently a proper ideal
has a common zero over $K$, by placing it inside a maximal ideal in this
\emph{set} ring. This proves the weak Nullstellensatz, including the
empty-zero criterion.

Suppose now that $G$ vanishes on all common zeros of the $F_i$ in $K^d$.
In $K[X_1,\ldots,X_d,T]$, the polynomials
\[
 F_1,\ldots,F_s,\quad 1-TG
\]
have no common zero. The weak statement gives an identity
\[
 1=\sum_i A_i(X,T)F_i(X)+B(X,T)(1-TG(X)).
\]
If $G=0$, certificate \eqref{eq:nullcertificate} is immediate. Otherwise
substitute $T=G^{-1}$ in the localized polynomial ring and multiply by a
sufficiently high power $G^N$, choosing $N\ge1$. This gives
\eqref{eq:nullcertificate} with all $H_i\in K[X]$. The reverse implication
follows by evaluating the identity at a common zero in any field and
using that $G^N=0$ implies $G=0$.

The original hypothesis about surcomplex zeros implies the hypothesis
about $K$-zeros, so it supplies the certificate in $K$. Conversely that
same identity holds in every larger workspace and hence at every
surcomplex tuple. This proves both directions without using a
proper-class maximal-ideal argument.
\end{proof}

The statement says nothing about an effective coefficient oracle for an
arbitrary surreal constant. Once coefficients belong to a computably
presented field, ordinary finite algebraic algorithms may be available;
an arbitrary Hahn support need not admit such a presentation. Likewise,
we have not asserted Noetherianity of the analytic coefficient rings in
the supplied manuscripts. Already over a point, for nonzero divisible $\Gamma$, the nonnegative
valuation ring $\Ov$ has the non-finitely generated ideal $\mv$: a finite
set of positive generator valuations has a least member $\delta>0$, but
$t^{\delta/2}$ is not in the ideal they generate over $\Ov$. The positive
part is not an ideal of the full Hahn field, where positive monomials
are units. This distinction is essential.

\subsection{Local perturbations do not control additional distant roots}
A monic fixed-degree one-variable family has constant fiber length by
\cref{thm:finitemap}. A local positive-support perturbation of an analytic
system is a different hypothesis. Even in one variable,
\begin{equation}\label{eq:escapingroot}
 X^2+tX^3=X^2(1+tX)
\end{equation}
has a double root in the infinitesimal monad of zero and a further root at
$-t^{-1}$. Its local reduction is $X^2$, but its global degree is three.
This is consistent with both the exact ball count and the local finite
intersection theorem of \cite{SourceResearch}. It is not consistent with
a claim that arbitrary positive-support polynomial perturbations preserve
all global roots without a degree or infinity condition. No such claim
is used in this article.

\section{Connection with the supplied global and analytic theories}\label{sec:connections}

\subsection{Why finite and infinite interpolation behave differently}
A finite union of well-ordered supports is well ordered. Therefore every
finite collection of root data and Hermite jets fits a Hahn workspace,
and \cref{thm:CRT} always gives a polynomial interpolant. There is no
extra infinite-union admissibility condition in this finite theorem.

The global manuscript \cite{SourceGlobal} treats a genuinely different
problem: infinitely many finite infinitesimal clusters above a closed
discrete ordinary subset of $\C$. Its exact divisor theorem tests the
union of supports of the \emph{symmetric cluster coefficients}, not the
supports of individually labelled roots. Its global interpolation theorem
tests the supports of polynomial remainder coordinates, not just raw
values at those roots. These source statements are not consequences of
finite CRT alone.

For example, the source proves that the simple roots $n+t^{1/n}$, one
for every positive integer $n$, cannot be the exact divisor of one
Hahn-coherent entire datum: the symmetric supports contain a descending
sequence. In contrast, all the roots of $(X-n)^n-t$ at every $n$ can be
realized, because the symmetric perturbation support is just $\{1\}$.
Selecting one root of each latter cluster by a value $1$ and the others
by $0$ has local remainder
\[
 B_n(U)=\frac1n\sum_{j=0}^{n-1}t^{-j/n}U^j.
\]
Its coefficient supports contain $-1+1/n$ and fail the required global
well-order condition. These examples and the nonexistence conclusions
are credited to \cite{SourceGlobal}. Here they identify exactly where a
finite algebraic theorem stops being an unrestricted infinite analytic
one.

\subsection{Polynomiality and coherence at every surreal scale}
The source \cite{SourceAnalysis} defines an all-scale Hahn-entire class
function by requiring, for every positive surreal $r$, an ordinary-domain
Hahn-coherent chart $F_r\in\mathcal H(B(0,2))$ satisfying
\[
 f(rw)=F_r\sh(w)\qquad(w\in B(0,2)\sh).
\]
It proves that these functions are precisely finite surcomplex
polynomials. We record its support mechanism because it explains the
natural role of the present subject.

If $a_n=f^{(n)}(0)/n!$ were nonzero for infinitely many $n$, the set of
valuations $v(a_n)$ would be bounded in $\No$. Choose a positive surreal
$\Delta$ larger than four times their absolute values, and use
$r=t^{-\Delta}$. The valuations
\[
 v(a_nr^n)=v(a_n)-n\Delta
\]
strictly decrease along successive nonzero Taylor indices. But the
coefficients of every derivative of one chart use only that chart's
well-ordered Hahn support. This is a contradiction. Thus the Taylor
series is a polynomial, and the source's coherent identity theorem
identifies it with $f$ on every chart. Conversely a finite polynomial
produces such charts at all scales by expanding its finitely many
coefficients into normal form.

This argument is inherited from the source, not claimed as new here. It
uses the full class of available scales and is not a classification of
entire functions in one fixed set-sized Hahn field. Our algebraic results
supply the finite polynomial calculus inside that rigid global category,
without requiring the all-scale hypothesis for each individual
polynomial statement.

\subsection{Compatibility, rather than conflation, of residues}
For a monic polynomial whose roots lie inside an appropriate ordinary
coherent contour, the source's coefficientwise residue functional equals
$[X^{n-1}]\remp_P h$ on polynomial numerators. This agreement is explicitly
proved in \cite{SourceResearch}. The algebraic functional in
\cref{sec:residue} is defined at every surcomplex scale and over any
coefficient ring, without a contour; on the shared domain its finite
partial-fraction formula gives the same local residue weights and jets.

Similarly, the parameter factor lifting of \cref{thm:parameterhensel}
preserves the source's fixed-domain and common-support requirements.
It does not replace the several-variable analytic division theorem by
ordinary polynomial division when a numerator or defining equation has
unbounded Taylor degree. The distinction is part of the construction,
not an afterthought.

\section{Scope, proof provenance, and further directions}\label{sec:scope}

\subsection{A contribution-by-contribution assessment}
\begin{longtable}{@{}P{0.30\textwidth}P{0.65\textwidth}@{}}
\toprule
Result or construction & Mathematical status and source of the proof\\
\midrule
\endfirsthead
\toprule Result or construction & Mathematical status and source of the proof\\
\midrule
\endhead
Finite factorization and CRT & Classical field algebra, proved here in set-sized Hahn workspaces; not a new algebraic-closure theorem.\\[0.5em]
Gauss--Lucas and root bounds & Classical finite inequalities and barycentric identities extended directly to real closed fields.\\[0.5em]
Order-modulus Rouch\'e & A precisely encoded fixed-degree transfer of the ordinary theorem using real-closed-field completeness. It is not a fine-topological contour theorem.\\[0.5em]
Exact coefficient uncertainty loci & Direct finite proofs of modulus and valuation versions, placed against the established pseudozero-set background. No priority claim for weighted root neighborhoods.\\[0.5em]
All-scale reduction and Newton polygons & Classical valued-field mechanisms given an explicit arbitrary-rank, arbitrary-center surcomplex formulation, including exact residue directions.\\[0.5em]
Critical-point allocation & A direct characteristic-zero reduction argument; related to the established non-Archimedean Rolle tradition. The branching formula is derived explicitly.\\[0.5em]
Cluster stability and precision bounds & Derived quantitative matching theorem, preservation of every pairwise valuation, a discriminant-only sufficient bound, and a strict-boundary collision example. No exhaustive novelty claim is made.\\[0.5em]
Support-controlled factor lifting & Coprime Hensel recursion with explicit positive support monoid, plus fixed-domain holomorphic parameters; related support methods are credited to the supplied papers.\\[0.5em]
Universal residue duality & Classical monogenic algebra developed over any commutative ring, with explicit dual basis and trace element. Source residue examples are identified as inherited.\\[0.5em]
Finite maps and Nullstellensatz & Classical algebraic results with full finite-data and class-size conventions; not claims about arbitrary analytic properness or Noetherianity.\\
\bottomrule
\end{longtable}

The literature check used the primary papers and author or publisher
records listed in the bibliography, consulted on September 21, 2026.
The search did not attempt to establish that every displayed extension
is absent under equivalent valued-field terminology. In particular,
``surcomplex'' does not by itself make a theorem new when the proof is
valid over every algebraically closed valued field of residue
characteristic zero.

\subsection{What the proofs do not assume}
All degrees are finite. Every invocation of an algebraically closed field,
maximal ideal, or finite-dimensional algebra takes place in a set-sized
workspace. Every factor-lifting recursion is supported on a well-ordered
positive monoid and has finite coefficient dependencies. No argument
assumes that ordinary sequential convergence detects fine surreal
convergence, that repeated addition of one positive value is cofinal, or
that a real-parameter path traces a nonconstant fine-continuous circle.

The stability theorem concerns squarefree polynomials and becomes singular
at a collision for a reason: individual roots then cease to be a stable
coordinate system. The residue pairing, by contrast, is deliberately
formulated on the full finite quotient algebra and remains perfect there.
These are compatible conclusions about different invariants.

The article proves no general effective procedure for arbitrary surreal
input coefficients. Finite algebraic identities and finite dependency at
one Hahn exponent do not make an arbitrary support computably enumerable.
The included symbolic computations verify selected finite formulas, not
these representability or summability statements in full generality.

\subsection{Concrete directions beyond this article}
A natural continuation is a multivariable analogue of the quantitative
cluster-stability theorem using finite local algebras and a controlled
Jacobian or different ideal, rather than pairwise root distances.
The supplied several-variable framework provides finite free algebras
for certain coupled systems, but its analytic support conditions must be
retained when introducing perturbation invariants.

Another direction is a relative description of root clusters over
parameter strata where the discriminant valuation changes. The present
factor-lifting and residue theorems remain coherent on their stated
domains, but a global labelled cluster tree can fail in the presence of
monodromy or collisions. These questions call for a carefully specified
parameter category, not a topological analogy with ordinary complex
coverings.

\subsection{Conclusion}
Finite surcomplex polynomial algebra is unusually robust: it retains
classical algebraic identities, real-closed geometric inequalities, and
precise non-Archimedean information simultaneously. The normalization
$P(a+t^\delta U)$ makes the different scales visible through one ordinary
reduction polynomial. Differentiation then counts critical points,
discriminants aggregate finite cluster depths, and a quantitative
precision threshold preserves the entire pairwise valuation geometry.
On the algebraic side, the residue functional and its explicit dual basis
remain valid even when individual root descriptions collide. These
results form a finite, all-scale polynomial counterpart to the analytic
and global support theories in the supplied manuscripts.

\appendix
\section{Notation and hypothesis reference}\label{app:notation}
\begin{longtable}{@{}P{0.28\textwidth}P{0.67\textwidth}@{}}
\toprule
Notation & Meaning\\
\midrule
\endhead
$\SC=\No[i]$ & Surcomplex class field; every polynomial has ordinary finite degree.\\[0.3em]
$t^\gamma=\omega^{-\gamma}$ & Conway monomial; positive exponent means an infinitesimal scale.\\[0.3em]
$K_\Gamma$ & Set-sized Hahn field $\C((t^\Gamma))$; algebraically closed when $\Gamma$ is divisible.\\[0.3em]
$v,\lc$ & Least support exponent and its nonzero complex coefficient.\\[0.3em]
$\Ov,\mv,\red$ & Finite valuation ring, infinitesimal ideal, and reduction to $\C$.\\[0.3em]
$\Bge{a}{\delta},\Bgt{a}{\delta}$ & Closed and open valuation balls, not modulus disks.\\[0.3em]
$w_{a,\delta}(P)$ & Minimum $\min_j(v(p_j)+j\delta)$ for $P(a+Z)=\sum p_jZ^j$.\\[0.3em]
$R_{a,\delta}$ & Nonzero reduction of $t^{-w_{a,\delta}(P)}P(a+t^\delta U)$.\\[0.3em]
$E^*$ & Additive monoid generated by a well-ordered positive support, including zero.\\[0.3em]
$\HH(U)$ & Hahn-coherent coefficient data $\OO(U)((t^\Gamma))$ on a fixed ordinary domain.\\[0.3em]
$d_{ij},\sigma_i,\delta_i$ & Pairwise distance valuations, derivative valuations, and closest-root scales in \eqref{eq:stabilitydata}.\\[0.3em]
$T(P),D$ & Root-specific precision threshold and discriminant valuation, for monic squarefree polynomials with finite roots.\\[0.3em]
$\lambda_P$ & Highest remainder coefficient, equivalently the formal $X^{-1}$ coefficient of $h/P$.\\[0.3em]
$b^i,G,M_h$ & Residue-dual basis, residue Gram matrix, and multiplication matrix of $h$.\\
\bottomrule
\end{longtable}

\section{Support and proof-dependency audit}\label{app:support}
\begin{longtable}{@{}P{0.36\textwidth}P{0.59\textwidth}@{}}
\toprule
Operation & Why it is defined, and what controls it\\
\midrule
\endhead
Polynomial evaluation and differentiation & Finitely many field additions and multiplications; the coefficient field is held constant under differentiation in $X$.\\[0.4em]
Finite CRT and Hermite interpolation & Inverses of finitely many nonzero root separations or coprime-factor classes; finite Taylor jets only.\\[0.4em]
Affine scaling and reduction & Finitely many coefficient operations, followed by the minimum of finitely many valuations and coefficient-zero reduction.\\[0.4em]
Coprime factor lifting & A recursion on $E^*$; nonlinear terms have strictly earlier positive exponents, and each output receives finitely many support-word contributions.\\[0.4em]
Parameter factor lifting & The same recursion; one finite holomorphic inverse matrix on the fixed ordinary domain introduces no new exponent.\\[0.4em]
Coherent parameter evaluation & Input support plus the positive monoid generated by finitely many displacement supports, as in \eqref{eq:coherentevaluation}.\\[0.4em]
Laurent expansion at infinity & Formal expansion in a separate indeterminate $X^{-1}$, with finitely determined coefficients; not an infinite evaluation at a surreal point.\\[0.4em]
Residue Gram matrix and dual basis & Finite polynomial operations in the coefficients of a monic polynomial; the determinant is a constant unit.\\[0.4em]
Root matching and cluster sums & Finite ball counts and finite ultrametric comparisons; no iterative convergence assertion.\\[0.4em]
Nullstellensatz certificate & Ordinary set-sized polynomial rings and a finite identity, which remains true under scalar extension.\\
\bottomrule
\end{longtable}

The geometric dependency chain is
\[
 \begin{gathered}
 \text{Hahn workspace and finite factorization}
 \ \Longrightarrow\ \text{weighted reduction and exact root counts}\\
 \Longrightarrow\ \text{critical-ball counts and valuation Rouch\'e}
 \ \Longrightarrow\ \text{root matching and discriminant precision}.
 \end{gathered}
\]
The residue chain is independent of the stability theorem:
\[
 \begin{gathered}
 \text{monic division}
 \Longrightarrow\text{explicit residue-dual basis}\\
 \Longrightarrow\text{trace element }P'
 \Longrightarrow\text{trace determinant and collision comparison}.
 \end{gathered}
\]
Order-modulus Rouch\'e uses a separate classical input, real-closed-field
transfer. The Hensel recursion uses the support lemma and a finite
coprime-factor splitting, not any of the fine-topological convergence
properties rejected in the source.

\section{Reproducible finite checks}\label{app:checks}
The archive contains \texttt{verify\_examples.py}, an exact-arithmetic
SymPy script, and its generated \texttt{verification.txt}. It checks
finite polynomial identities and selected leading terms, including the
Newton example, the four-root discriminant and derivative weights,
coprime factor lifting, the residue dual basis and its Gram determinant,
the trace formula, and the branch-value discriminant identity.
The report records the actual ranges and cases checked.

These computations are intentionally not described as a formal
verification of the article. They cannot prove arbitrary-group
well-ordering, a transfinite recursion, real-closed transfer, or the
universal quantified stability statement by finite sampling. They check
algebraic formulas that are particularly prone to sign or coefficient
errors. The general assertions rely on the mathematical proofs above.

The source is standalone, with an internal bibliography. Compile with
\begin{center}\small
\texttt{latexmk -pdf surcomplex\_polynomial\_algebra.tex}
\end{center}
or run \texttt{pdflatex} until cross-references stabilize. No external image,
bibliography database, or network resource is needed to compile it.

\begin{thebibliography}{99}
\addcontentsline{toc}{section}{References}
\small
\bibitem{SourceAnalysis}
\emph{Surcomplex Analysis: Infinitesimal Calculus, Hahn-Coherent Holomorphy,
and Contour Theory}. Supplied manuscript, September 21, 2026,
\texttt{surcomplex\_analysis(3).tex}. Source of the three-level analytic
framework, strong evaluation, one-variable preparation and residues, and
all-scale rigidity. No publication or personal authorship is inferred
from the filename.

\bibitem{SourceResearch}
\emph{Hartogs Coherence, Finite Maps, and Residue Duality over the
Surcomplex Numbers}. Supplied manuscript, September 21, 2026,
\texttt{surcomplex\_research.tex}. Source of the fixed-domain parameter
and finite-intersection framework, one-variable coherent residue
identities, and the credited cubic collision example.

\bibitem{SourceGlobal}
\emph{Global Divisors and Cohomological Obstructions in Hahn-Coherent
Surcomplex Analysis}. Supplied manuscript, September 21, 2026,
\texttt{surcomplex\_global\_theorems.tex}. Source of the infinite symmetric
support criterion, restricted interpolation targets, and global
nonexistence examples discussed in \cref{sec:connections}.

\bibitem{Gonshor}
Harry Gonshor, \emph{An Introduction to the Theory of Surreal Numbers},
London Mathematical Society Lecture Note Series \textbf{110},
Cambridge University Press, 1986.

\bibitem{Neumann}
B. H. Neumann, ``On ordered division rings,''
\emph{Transactions of the American Mathematical Society}
\textbf{66} (1949), 202--252.

\bibitem{Poonen}
Bjorn Poonen, ``Maximally complete fields,''
\emph{L'Enseignement Math\'ematique} (2) \textbf{39} (1993), 87--106.
\href{https://math.mit.edu/~poonen/papers/amsval.pdf}{Author-hosted full text}.
Corollary~4 supplies the algebraic-closure input for divisible Hahn
workspaces.

\bibitem{Tarski}
Lou van den Dries, ``Alfred Tarski's elimination theory for real closed
fields,'' \emph{The Journal of Symbolic Logic} \textbf{53} (1988),
no.~1, 7--19.
\href{https://doi.org/10.2307/2274424}{doi:10.2307/2274424}.

\bibitem{CohenMahboubi}
Cyril Cohen and Assia Mahboubi,
``Formal proofs in real algebraic geometry: from ordered fields to
quantifier elimination,'' \emph{Logical Methods in Computer Science}
\textbf{8} (2012), no.~1, article~2, 1--40.
\href{https://lmcs.episciences.org/844}{Publisher record};
\href{https://arxiv.org/abs/1201.3731}{arXiv:1201.3731}.
This reference concerns the classical logical input, not a formal
verification of the present manuscript.

\bibitem{Eisermann}
Michael Eisermann, ``The Fundamental Theorem of Algebra made effective:
an elementary real-algebraic proof via Sturm chains,''
\href{https://arxiv.org/abs/0808.0097}{arXiv:0808.0097}, version~4, 2012.
Used for comparison with real-algebraic polynomial root counting,
not as a new surcomplex algorithm claim.

\bibitem{Eberl}
Manuel Eberl, ``Two theorems about the geometry of the critical points of
a complex polynomial,'' \emph{Archive of Formal Proofs}, November 14,
2023.
\href{https://isa-afp.org/entries/Polynomial_Crit_Geometry.html}{Formal development record}.
Ordinary complex Gauss--Lucas and Jensen theorems.

\bibitem{FaroukiHan}
Rida T. Farouki and Chang Yong Han,
``Robust plotting of generalized lemniscates,''
\emph{Applied Numerical Mathematics} \textbf{51} (2004), nos.~2--3,
257--272.
\href{https://doi.org/10.1016/j.apnum.2004.05.007}{doi:10.1016/j.apnum.2004.05.007}.
The publisher abstract describes weighted pseudozero sets and their
root-neighborhood interpretation; no detailed numerical algorithm from
this paper is used here.

\bibitem{Faber}
Xander Faber, ``Topology and Geometry of the Berkovich Ramification Locus
for Rational Functions, II,''
\href{https://arxiv.org/abs/1104.0943}{arXiv:1104.0943}, version~3, 2013.
\href{https://doi.org/10.1007/s00208-012-0872-3}{doi:10.1007/s00208-012-0872-3}.
Related non-Archimedean ramification and Rolle theory. No assumption is
made that a rank-one analytic construction automatically applies to every
surreal value group.

\bibitem{StacksNS}
The Stacks Project Authors, \emph{The Stacks Project},
\href{https://stacks.math.columbia.edu/tag/00FV}{Tag~00FV: Hilbert
Nullstellensatz}. Consulted September 21, 2026.

\end{thebibliography}
\end{document}
